% Options for packages loaded elsewhere
\PassOptionsToPackage{unicode}{hyperref}
\PassOptionsToPackage{hyphens}{url}
\documentclass[
  10pt,
]{article}
\usepackage{xcolor}
\usepackage[margin=20mm]{geometry}
\usepackage{amsmath,amssymb}
\setcounter{secnumdepth}{-\maxdimen} % remove section numbering
\usepackage{iftex}
\ifPDFTeX
  \usepackage[T1]{fontenc}
  \usepackage[utf8]{inputenc}
  \usepackage{textcomp} % provide euro and other symbols
\else % if luatex or xetex
  \usepackage{unicode-math} % this also loads fontspec
  \defaultfontfeatures{Scale=MatchLowercase}
  \defaultfontfeatures[\rmfamily]{Ligatures=TeX,Scale=1}
\fi
\usepackage{lmodern}
\ifPDFTeX\else
  % xetex/luatex font selection
  \setmainfont[BoldFont=lmroman10-bold.otf,ItalicFont=lmroman10-italic.otf,BoldItalicFont=lmroman10-bolditalic.otf]{lmroman10-regular.otf}
  \setmonofont[BoldFont=lmmonolt10-bold.otf,ItalicFont=lmmono10-italic.otf]{lmmono10-regular.otf}
  \setmathfont[]{latinmodern-math.otf}
\fi
% Use upquote if available, for straight quotes in verbatim environments
\IfFileExists{upquote.sty}{\usepackage{upquote}}{}
\IfFileExists{microtype.sty}{% use microtype if available
  \usepackage[]{microtype}
  \UseMicrotypeSet[protrusion]{basicmath} % disable protrusion for tt fonts
}{}
\makeatletter
\@ifundefined{KOMAClassName}{% if non-KOMA class
  \IfFileExists{parskip.sty}{%
    \usepackage{parskip}
  }{% else
    \setlength{\parindent}{0pt}
    \setlength{\parskip}{6pt plus 2pt minus 1pt}}
}{% if KOMA class
  \KOMAoptions{parskip=half}}
\makeatother
\usepackage{longtable,booktabs,array}
\newcounter{none} % for unnumbered tables
\usepackage{calc} % for calculating minipage widths
% Correct order of tables after \paragraph or \subparagraph
\usepackage{etoolbox}
\makeatletter
\patchcmd\longtable{\par}{\if@noskipsec\mbox{}\fi\par}{}{}
\makeatother
% Allow footnotes in longtable head/foot
\IfFileExists{footnotehyper.sty}{\usepackage{footnotehyper}}{\usepackage{footnote}}
\makesavenoteenv{longtable}
\usepackage{graphicx}
\makeatletter
\newsavebox\pandoc@box
\newcommand*\pandocbounded[1]{% scales image to fit in text height/width
  \sbox\pandoc@box{#1}%
  \Gscale@div\@tempa{\textheight}{\dimexpr\ht\pandoc@box+\dp\pandoc@box\relax}%
  \Gscale@div\@tempb{\linewidth}{\wd\pandoc@box}%
  \ifdim\@tempb\p@<\@tempa\p@\let\@tempa\@tempb\fi% select the smaller of both
  \ifdim\@tempa\p@<\p@\scalebox{\@tempa}{\usebox\pandoc@box}%
  \else\usebox{\pandoc@box}%
  \fi%
}
% Set default figure placement to htbp
\def\fps@figure{htbp}
\makeatother
\setlength{\emergencystretch}{3em} % prevent overfull lines
\providecommand{\tightlist}{%
  \setlength{\itemsep}{0pt}\setlength{\parskip}{0pt}}
\usepackage{newunicodechar}
\newunicodechar{β}{\ensuremath{\beta}}
\newunicodechar{γ}{\ensuremath{\gamma}}
\newunicodechar{δ}{\ensuremath{\delta}}
\newunicodechar{η}{\ensuremath{\eta}}
\newunicodechar{θ}{\ensuremath{\theta}}
\newunicodechar{κ}{\ensuremath{\kappa}}
\newunicodechar{μ}{\ensuremath{\mu}}
\newunicodechar{π}{\ensuremath{\pi}}
\newunicodechar{σ}{\ensuremath{\sigma}}
\newunicodechar{χ}{\ensuremath{\chi}}
\newunicodechar{Δ}{\ensuremath{\Delta}}
\newunicodechar{ᵢ}{\ensuremath{{}_{i}}}
\newunicodechar{⁴}{\ensuremath{{}^{4}}}
\newunicodechar{⁻}{\ensuremath{{}^{-}}}
\newunicodechar{₀}{\ensuremath{{}_{0}}}
\newunicodechar{₃}{\ensuremath{{}_{3}}}
\newunicodechar{₄}{\ensuremath{{}_{4}}}
\newunicodechar{→}{\ensuremath{\rightarrow}}
\newunicodechar{↔}{\ensuremath{\leftrightarrow}}
\newunicodechar{∈}{\ensuremath{\in}}
\newunicodechar{≈}{\ensuremath{\approx}}
\newunicodechar{≥}{\ensuremath{\geq}}
\newunicodechar{≳}{\ensuremath{\gtrsim}}
\newunicodechar{−}{\ensuremath{-}}
\BeforeBeginEnvironment{longtable}{\begingroup\footnotesize}
\AfterEndEnvironment{longtable}{\endgroup}
\setlength{\emergencystretch}{5em}
\newunicodechar{≡}{\ensuremath{\equiv}}
\usepackage{bookmark}
\IfFileExists{xurl.sty}{\usepackage{xurl}}{} % add URL line breaks if available
\urlstyle{same}
\hypersetup{
  hidelinks,
  pdfcreator={LaTeX via pandoc}}

\author{}
\date{}

\begin{document}

\section{Can a photonic state-space model be trained on-chip? A
pre-registered in-situ-training bake-off on a realistic silicon-nitride
ring
substrate}\label{can-a-photonic-state-space-model-be-trained-on-chip-a-pre-registered-in-situ-training-bake-off-on-a-realistic-silicon-nitride-ring-substrate}

\emph{Lucas Talandier --- independent researcher, Paris}

\textbf{Abstract.} To our knowledge, no physical photonic system has yet
had the parameters that define a continuous-time dissipative-resonator
recurrence --- pole positions and inter-resonator couplings, the physics
that \emph{is} the memory --- trained on the device by gradient-based or
gradient-estimating methods.¹ We ask whether such training is feasible
for a photonic state-space model on ultra-low-loss silicon nitride, and
answer it in simulation: \textbf{on a pre-registered dissipative
substrate model, both hardware-committed methods train the recurrence to
within margin of the exact-gradient ceiling on 8/8 seeds at realistic
noise.} Every threshold was frozen before the run that consumed it.

We derive the SSM↔ring mapping and its realizable pole region, build one
shared substrate (finite Q, saturating gain, amplifier noise), and run a
four-method bake-off. PAT needs 4.6× fewer device passes than model-free
SPSA. The partial energy budget assigns SPSA 2.6 mJ in conversion energy
and PAT 13--44 J in digital-twin energy; total training energy remains
unmeasured. Hamiltonian-echo learning is censored --- a quantified
feasibility bound, the substrate's own dissipation defeating the echo. A
measured controllability profile (one drive trains ≈3 of 32 rings; four
taps recover all 32) makes the input map a first-class design axis ---
and the winning routes' trained solutions preserve it, converging on an
estimator-independent damp-the-driven-rings profile whose interior
contribution survives a registered taps-only falsifier, narrowly.

Against the decisive baseline --- calibrate offline, deploy, retrain the
readout --- in-situ training shows no resolved advantage at five
calibration-error levels spanning 5--30\% after a corrected
command-binding rerun, or under common-mode drift. Under uncorrelated
per-ring drift it holds a pre-registered 2.42× advantage (a
threshold-crossing under a frozen rule; the ratio's own CI spans {[}1.7,
4.6{]}). A registered inline-envelope follow-up finds no
energy-advantage window at 0.1--2 GS/s against a per-tap digital
equalizer (best digital/photonic ratio 0.64). Capacity remains a
separate limitation: a registered ablation finds the deployed
equalizer's output rides on 6--8 of 32 rings, and no workload in our
data exercises more. One registered prediction --- that the damping
optimum tracks task memory span --- failed, and is reported as failed.
The pre-registration ledger, substrate model, and training code are
released with the paper.

¹ \emph{The nearest neighbor, an in-situ-trained optical recurrent
network, trains interferometer weights around an optoelectronic relay;
its resonators stay fixed.}

\subsection{1. Introduction}\label{introduction}

\subsubsection{1.1 The gap}\label{the-gap}

Photonic neural networks have learned to train themselves. In
feedforward meshes and diffractive processors, on-chip training is now
routine enough to have families: model-free perturbative methods {[}1,
2{]}, hybrid physical-forward/digital-backward methods {[}3{]}, and
in-situ adjoint methods that read gradients from interference {[}4{]}.
But the systems that most need on-device training --- \emph{recurrent}
photonic processors, whose memory lives in the physics --- have barely
begun to receive it. Photonic reservoir computing deliberately avoids
the problem: the recurrence is fixed, random, and only a readout is
trained {[}5{]}. Where internal parameters of a photonic recurrence have
been adjusted at all, it has been by calibration or regime-tuning rather
than task-driven training {[}6{]}, by reinforcement-style search over a
readout while the loop stays fixed {[}7{]}, or in systems whose
recurrent state is carried digitally between optical passes {[}8{]}. A
two-modality literature search with a pre-registered kill-criterion (§8;
supplementary) found no demonstration, by any method, of the following:

\begin{quote}
\textbf{a continuous-time dissipative-resonator recurrence --- pole
positions and inter-resonator couplings --- trained in situ on the
physical device by gradient-based or gradient-estimating methods on a
computational task.}
\end{quote}

That sentence is this program's target, with each qualifier
load-bearing: \emph{on a computational task} excludes the
servo/calibration lineage; \emph{physical parameters of the recurrence}
excludes hybrid-digital state carriage; \emph{weight-tied recurrence}
excludes feedforward meshes folded in time. Refreshes of the search at
assembly (2026-07-12) and a page-level verification round (2026-07-27),
followed by primary-source checks and a bounded search on 2026-09-13
(memos in supplementary), map the boundary against the strongest
2025--26 neighbors, which we dispatch by name because each is the
``nearest miss'' along one qualifier --- and one of them moved the
boundary. The \textbf{nearest neighbor} is the monolithic optical
recurrent accelerator of Wu et al.~{[}9{]}: its ORNN chip \emph{is}
trained in situ, by a model-free perturbative method (SPGD, two physical
evaluations per update) on a classification task, with a weight-tied
mesh applied across wavelength-encoded time steps --- to our knowledge
the first in-situ-trained optical recurrent network of any kind, and we
cite it as such. What it does not do is train the parameters in the
boxed sentence: its trained weights are interferometer-mesh voltages,
its recurrent state is re-generated electronically at every step through
a photodetector--modulator relay, and its resonators are calibrated once
and held static --- the continuous-time dissipative-resonator
recurrence, whose \emph{poles and couplings are themselves the memory},
remains untrained. The remaining near-misses each fall along one
qualifier: on-chip all-photonic backpropagation is now demonstrated ---
for a \emph{feedforward} network {[}10{]}; microring weight banks have
been trained in situ through on-chip optical backprop --- as
\emph{feedforward} layers {[}11{]}; a time-synthetic fiber-loop network
trains in situ --- with per-step distinct programmed parameters,
i.e.~unrolled feedforward rather than a weight-tied recurrence {[}12{]};
a modulation-and-weighting microring array trains weights sitting
\emph{inside} an analog recurrent loop on-device --- by particle-swarm
search, explicitly without gradients {[}13{]}; an optoelectronic delay
reservoir has had recurrence-defining parameters optimized in situ ---
by Bayesian search, through a digital feedback loop {[}14{]}; and a
silicon photonic reservoir equalizer is ``trained in hardware'' --- its
readout only, by an evolution strategy, the reservoir couplings fixed by
fabrication {[}15{]}. No coupled-resonator lattice has had its couplings
learned on-device by any method. (The claim is time-indexed: the search
refresh re-runs once more immediately before submission, with the
registered page-level reads listed in supplementary.)

\subsubsection{1.2 Why a state-space model, and why silicon
nitride}\label{why-a-state-space-model-and-why-silicon-nitride}

Two developments make this the right moment to close the gap. On the
algorithmic side, deep state-space models showed that \emph{linear}
recurrences with well-placed poles --- not gated nonlinear dynamics ---
carry most of long-sequence performance {[}16, 17{]}, and the
oscillatory LinOSS line extended this to second-order units that are
exactly coupled damped oscillators {[}18, 19{]}. A lattice of coupled
microrings \emph{is} such a system: poles are ring detunings and losses,
couplings are physical couplers, and --- decisively for hardware --- the
LinOSS stability analysis tolerates nonnegative damping, so a
lossy-but-high-Q dissipative lattice is in-regime rather than an
approximation to a conservative ideal (§2). Damping becomes a design
knob, not an embarrassment.

On the platform side, the figure of merit for a dissipative ring
recurrence is memory per pass --- round-trip loss, hence intrinsic Q.
Ultra-low-loss silicon nitride is class-leading and foundry-accessible
(\(Q_i\) near \(10^7\) on multi-project-wafer runs {[}20{]}), thermally
stable enough for the slow thermo-optic actuation that
gradient-estimating training wants, and needs so little gain that the
injected amplifier noise stays low (§3). The same choice has a cost we
do not hide: SiN's weak \(\chi^{(3)}\) makes the one training route that
needs an optical nonlinearity --- Hamiltonian-echo learning, whose echo
is phase conjugation --- \emph{harder}, not easier, and §4--§5 treat
that honestly rather than idealizing it away.

\subsubsection{1.3 What this paper does}\label{what-this-paper-does}

This is a theory-and-simulation (Stage-0) paper, built to be falsified
before it is fabricated. Its results are gated by a pre-registration
ledger --- margins, budgets, fairness conventions, and statistical rules
frozen (and PI-signed) before the runs that consume them, with the
commit trail shipped as supplementary material (§3.7). Concretely:

\begin{enumerate}
\def\labelenumi{\arabic{enumi}.}
\tightlist
\item
  \textbf{A mapping with a realizable region} (§2): the oscillatory-SSM
  ↔ coupled-SiN-ring correspondence, with the pole region bounded by
  measured platform numbers, an actuation map for \{δ, κ\_ext, μ\}, and
  the backscatter doublet carried as the model rather than a correction.
\item
  \textbf{A shared dissipative substrate} (§3): saturating Er gain at a
  registered operating point, Langevin ASE at NF 7 dB, three cells from
  foundry-floor to aspirational, a measured feasible box bounded by the
  lasing crossing, and a measured minimal multi-point input map after
  single-drive controllability was found to collapse (effective
  dimension ≈3 of 32 rings).
\item
  \textbf{A four-method bake-off under one fairness contract} (§4--§5):
  SPSA, PAT (with registered twin-mismatch), a recurrent in-situ adjoint
  (charged as if realizable; optimistic bound, labelled), and RHEL over
  a concrete χ³-FWM echo sub-model --- scored by device passes to a
  pre-registered target relative to the exact-gradient ceiling, eight
  seeds, censoring-aware statistics.
\item
  \textbf{The result} (§5): both hardware-committed methods train the
  recurrence to within margin of the ceiling on all eight seeds --- the
  sentence in §1.1, demonstrated in simulation at a foundry-class noise
  cell. The ranking (PAT \textless{} adjoint \textless{} SPSA on device
  passes; RHEL censored, and worse than readout-only under an honest
  echo) settles the hardware roadmap on PAT/SPSA without promotion of
  the exotic routes. And the comparison the fair design was built to
  expose shows no resolved calibration difference at five levels
  spanning 5--30\%-class error after the corrected binding rerun (§5.5),
  and no advantage under common-mode drift, which a laser re-lock
  absorbs --- while under uncorrelated per-ring drift, which a global
  laser re-lock cannot absorb, in-situ retraining holds a
  \textbf{declared, pre-registered 2.42× advantage} (both conditions of
  the frozen rule hold --- point ratio ≥ 2× \emph{and} difference-CI
  excluding zero; the coarse-floor estimate 1.84× is co-reported), a gap
  that grows with accumulated drift (§5.5, §7).
\item
  \textbf{Limits, stated as limits} (§8): robustness is measured against
  modelled imperfections only; the anchor risks, verification debts, and
  the single-session review period are disclosed with the same
  specificity as the results.
\end{enumerate}

Our contribution is a simulation test of training the pole positions and
couplings that define a dissipative SiN memory lattice, under registered
thresholds and explicit cost accounting. PAT and SPSA have prior chip
demonstrations; the literature distinction is the specific recurrence in
§1.1. A physical demonstration of this architecture remains future work.
Its systems value is assessed separately in §7, including the negative
inline follow-up.

\subsection{2. From oscillatory state-space models to coupled SiN
microrings}\label{from-oscillatory-state-space-models-to-coupled-sin-microrings}

\subsubsection{2.1 The model class}\label{the-model-class}

A structured state-space model (SSM) processes a sequence
\(u_1, u_2, \dots\) through a linear recurrence
\(x_{t+1} = \bar A x_t + \bar B u_t\),
\(y_t = \mathrm{Re}(\bar C x_t) + D u_t\), whose expressive power is set
by where the eigenvalues (poles) of \(\bar A\) can be placed and how
precisely. The diagonal variants that dominate current practice --- S4D
and DSS {[}17, 21{]} --- reduce \(\bar A\) to a set of independent
complex poles, \(H(s) = \sum_j c_j b_j/(s - \lambda_j) + D\); the
oscillatory LinOSS family {[}18{]} instead parameterizes forced, damped
harmonic oscillators, and its damped extension D-LinOSS {[}19{]} makes
the damping of each mode a trainable parameter. Crucially for what
follows, the LinOSS construction is valid for any nonnegative-diagonal
(dissipative) state matrix: stability does not require conservative
dynamics, only that every mode decay. A physical substrate that is
\emph{lossy but slowly so} is therefore in-regime by construction ---
the loss is not an error term to be fought but the damping parameter of
the model class itself. (Where in the damping range the trained optimum
actually lands --- and what that does to the LinOSS-specific part of
this motivation --- is measured, not assumed; see §6.)

\subsubsection{2.2 One ring is one trainable complex
pole}\label{one-ring-is-one-trainable-complex-pole}

The mode amplitude \(a_j\) of a silicon-nitride microring obeys temporal
coupled-mode theory (Haus energy-amplitude convention, \(|a|^2\) =
stored energy) {[}22{]}:

\[\dot a_j = (i\delta_j - \kappa_{\mathrm{tot},j})\,a_j
+ i\sum_{k\neq j}\mu_{jk}\,a_k + \sqrt{2\kappa_{\mathrm{ext},j}}\,u(t),\]

with
\(\kappa_{\mathrm{tot},j} = \kappa_{i,j} + 2\kappa_{\mathrm{ext},j}\)
the amplitude decay rate (intrinsic + both bus ports of the add--drop
registry ring, each loading the mode at \(\kappa_\mathrm{ext}\); the
drive enters through one port, hence the single
\(\sqrt{2\kappa_{\mathrm{ext}}}\) injection term), \(\delta_j\) the
detuning of the drive from the ring resonance, and \(\mu_{jk}\) the
inter-ring coupling. A single uncoupled ring is therefore exactly one
\textbf{complex} Laplace pole,

\[s_j = -\kappa_{\mathrm{tot},j} + i\,\delta_j,\]

--- complex, not real, because the optical envelope carries a carrier. A
bank of \(N\) rings natively realizes a \textbf{diagonal complex-pole
SSM of the S4D/DSS class}: one ring per pole, with the injection and
readout couplings playing the role of the residues \(c_j b_j\). The
oscillatory LinOSS unit is recovered \emph{exactly} as the uncoupled,
real-input/output special case: a real damped second-order oscillator is
a conjugate pole pair, and the identification against the LinOSS
eigenvalue \(a_i = -e^{\alpha_i} + i\beta_i\) is

\[e^{\alpha_i} = \kappa_\mathrm{tot}, \qquad \beta_i = \delta.\]

Storing \(\alpha_i = \log\kappa_\mathrm{tot}\) enforces the dissipative
sign and \(\kappa_\mathrm{tot} > 0\) by construction --- the physical
substrate implements the model's own stability parameterization. We
deliberately claim the broader class and treat LinOSS as its special
case, with one consequence made explicit now: published LinOSS/D-LinOSS
benchmark results validate only the \(\mu = 0\) diagonal reduction, and
transfer to the coupled realization is verified in-house against a
BPTT-on-substrate ceiling (§5.1) rather than assumed.

Discretization is exact, not approximate: for piecewise-constant input
the zero-order-hold (van-Loan matrix-exponential) step gives discrete
poles \(z = e^{s\,dt}\) to machine precision, with \(dt\) the round-trip
time. The per-step memory retention is
\(|z| = e^{-\kappa_\mathrm{tot} dt}\); long memory is literally
\(|z| \to 1\).

\textbf{Validation.} The mapping is validated at three levels
(test-anchored; all thresholds pre-specified): (i) \emph{by
construction} --- uncoupled system poles match the CMT reference to
\(<10^{-3}\) and discrete \(|z|\) to \(<10^{-9}\); (ii) \emph{in the CW
limit} --- the dynamical model's steady state recovers independently
derived static transfer functions with error scaling as
\(O(1/\mathcal{F})\) (finesse), measured \(4.5\times10^{-4}\) at
\(\mathcal{F} = 1048\) falling to \(4.4\times10^{-5}\) at
\(\mathcal{F} = 10473\); the registry rings sit at
\(\mathcal{F} \approx 10^3\)-- \(1.5\times10^4\), where the CMT pole is
an excellent model; (iii) \emph{in the time domain} --- against an
independent RK45 integration, ringdown matches the closed form to
\(<10^{-9}\) and both pole coordinates (\(\kappa\) and \(\delta\))
re-fitted from the trajectory recover the set values to \(<10^{-5}\).
All three are \emph{internal} checks --- independent implementations and
limits of the same coupled-mode model class --- so they certify that the
simulator faithfully computes the model it claims, not that the model
captures a fabricated device; the latter is a hardware question, scoped
honestly in §8.1.

\subsubsection{2.3 Coupling rings: the structured off-diagonal
generalization}\label{coupling-rings-the-structured-off-diagonal-generalization}

Energy-conserving inter-ring coupling enters as \(i\Omega\) with
\(\Omega\) real-symmetric --- an anti-Hermitian addition to the state
matrix. It hybridizes the poles (photonic-molecule supermodes) while
conserving total damping: the eigenvalues move along the imaginary axis
while
\(\sum_j \mathrm{Re}\,\lambda_j = -\sum_j \kappa_{\mathrm{tot},j}\) is
invariant (verified: two-ring beat frequency equals the eigenvalue
splitting \(2\mu\) to \(<2\%\), with a no-beat control at \(\mu = 0\)).
The trainable \(\mu_{jk}\) are thus a \emph{mild, structured}
generalization beyond diagonal-\(A\) SSMs --- off-diagonal terms that
reshape frequencies but cannot buy stability or memory. They are also
load-bearing for the claim this paper builds toward: inter-ring
couplings are recurrence-defining parameters (§2.5).

\subsubsection{2.4 The realizable pole
region}\label{the-realizable-pole-region}

Where can the poles actually go? Four physical boundaries define the
envelope (Fig. F1).

\textbf{Memory is loss-limited; stability is free.} Passive rings
satisfy \(\kappa_\mathrm{tot} \geq \kappa_i = \omega_0/(2Q_i)\), so the
maximum amplitude (state) memory time is \(1/\kappa_i = 2Q_i/\omega_0\)
--- twice the photon-energy lifetime, a convention this paper uses
consistently. Across the registered platform range this gives
\textbf{3.29 ns (329 round trips)} at the foundry-conservative corner
(\(Q_i = 2\times10^6\)) to \textbf{49.4 ns (4937 round trips)} at the
class-leading corner (\(Q_i = 3\times10^7\)): a 15× memory span that is
the platform's dominant figure of merit and the reason this program is
on SiN (§1). Net optical gain moves \(\kappa_\mathrm{tot} \to 0\) and
hence \(|z| \to 1\), extending memory past the passive floor at the cost
of saturation and amplified-spontaneous-emission noise --- deferred to
the substrate model (§3).

\textbf{Frequency is FSR-bounded.} The detuning axis aliases at one free
spectral range: \(|\beta_j\,dt| \leq \pi\) with \(dt = 1/\mathrm{FSR}\).
The reachable imaginary axis is one FSR wide.

\textbf{Readout trades against memory.} The external coupling sets both
memory (\(1/\kappa_\mathrm{tot}\), maximized by deep undercoupling) and
readability (residue \(\propto 2\kappa_\mathrm{ext}\); on-resonance drop
efficiency \((2\kappa_\mathrm{ext}/\kappa_\mathrm{tot})^2\),
i.e.~detector SNR) in opposite directions; their product is bounded by
the passive memory length (tested). At the foundry corner the trade runs
from 274 round trips at drop efficiency 0.028 (deep undercoupling,
\(\kappa_\mathrm{ext} = 0.1\kappa_i\)) to 16 round trips at 0.91 (deep
overcoupling, \(10\kappa_i\)); the same Pareto shape holds at the
class-leading corner with the memory axis scaled \textasciitilde15×. The
operating point is not chosen here: it is pre-registered (§3, PR-4), and
\(\kappa_\mathrm{ext}\) is itself a pole-real-part actuator, so a
trainable-\(\kappa_\mathrm{ext}\) policy folds the readout trade into
the recurrence training.

\textbf{Backscatter bounds the one-ring-one-pole abstraction --- by
roughness, not cleanly by \(Q\).} Sidewall roughness couples the
counter-propagating modes at a coherent rate \(\gamma\) that is a
fabrication constant, essentially independent of
\(\kappa_\mathrm{tot}\); the resonance resolves into a standing-wave
doublet when \(2\gamma \gtrsim \kappa_\mathrm{tot}\) (conservative HWHM
criterion; the FWHM criterion shifts every crossover ×2, conclusion
unchanged) {[}23{]}. Because \(\kappa_i\) falls as \(Q_i\) rises while
\(\gamma\) does not, splitting \emph{grows} with \(Q\): from measured
SiN rates, a clean damascene-class process (\(\gamma/2\pi \approx 12\)
MHz {[}24{]}) keeps the single-pole picture at the foundry corner
(\(2\gamma/\kappa_i
\approx 0.5\)) and breaks it at \(Q \approx 4\times10^6\) (HWHM;
\(8\times10^6\) FWHM), while a rough subtractive process splits 21--75
\% of resonances \emph{already at} \(Q_i \lesssim 2.7\times10^6\), with
average doublet separations of 180--320 MHz depending on etch mask
{[}25{]} --- i.e. modal-coupling rates
\(\gamma/2\pi \approx 90\)--\(160\) MHz under the standard
\(2\gamma\)-separation convention {[}23{]}, a mapping we state here
because the source tabulates separations, not rates (page-verified
2026-07-12). \emph{The assembled crossover curve brackets published SiN
data points --- no single published SiN crossover exists, and no
\(\gamma\) is published for the specific target processes;
primary-source verification is registered before submission.} Two
consequences propagate forward: the substrate model carries a
roughness-gated CW/CCW doublet knob, default ON except at the
clean-damascene corner (§3); and a platform tension is on record --- the
best-memory (highest-\(Q\)) platforms are the most splitting-prone,
while the splitting-safe low-\(Q\) foundry corner is memory-poor.
Overcoupling suppresses the visible doublet, so the
\(\kappa_\mathrm{ext}\) policy and the splitting risk are coupled: the
max-memory (deep-undercoupled) regime is the most exposed.

\subsubsection{2.5 What ``training the recurrence in situ'' means
physically}\label{what-training-the-recurrence-in-situ-means-physically}

The recurrence is defined by the state matrix
\(M = \mathrm{diag}(-\kappa_{\mathrm{tot},j} + i\delta_j) + i\Omega\).
Every entry has a concrete, independently addressable actuator:

{\def\LTcaptype{none} % do not increment counter
\begin{longtable}[]{@{}
  >{\raggedright\arraybackslash}p{(\linewidth - 6\tabcolsep) * \real{0.2500}}
  >{\raggedright\arraybackslash}p{(\linewidth - 6\tabcolsep) * \real{0.2500}}
  >{\raggedright\arraybackslash}p{(\linewidth - 6\tabcolsep) * \real{0.2500}}
  >{\raggedright\arraybackslash}p{(\linewidth - 6\tabcolsep) * \real{0.2500}}@{}}
\toprule\noalign{}
\begin{minipage}[b]{\linewidth}\raggedright
Pole quantity
\end{minipage} & \begin{minipage}[b]{\linewidth}\raggedright
Physical knob
\end{minipage} & \begin{minipage}[b]{\linewidth}\raggedright
Actuator
\end{minipage} & \begin{minipage}[b]{\linewidth}\raggedright
Cost
\end{minipage} \\
\midrule\noalign{}
\endhead
\bottomrule\noalign{}
\endlastfoot
\(\mathrm{Im}\) (frequency) \(\delta_j\) & ring resonance offset &
thermo-optic heater (trench-isolated), one per ring & low: 1 heater +
DAC channel/ring \\
\(\mathrm{Re}\) (damping) \(\kappa_{\mathrm{tot},j}\) & bus--ring
coupling \(\kappa_{\mathrm{ext},j}\) & tunable coupler (MZI-assisted
gap) & medium: 1--2 heaters/ring \\
same, past the passive floor & per-ring net gain \(g_j\) & Er:SiN /
III--V pump current & high (gain stage + ASE cost); not required for the
claim \\
off-diagonal \(\mu_{jk}\) & ring--ring coupling & tunable
photonic-molecule coupler (or bus-mediated mesh) & medium--high;
topology fixed at fab, strengths trainable \\
residues \(B, C\) & injection/readout mesh & thermo-optic MZI mesh &
\textbf{not recurrence-defining} \\
\end{longtable}
}

The last row is the load-bearing exclusion. Training only the
injection/readout mesh --- the residues --- while the poles stay fixed
is precisely reservoir computing with a trained linear head, and prior
photonic work in that regime (including reinforcement-learning-tuned
readouts {[}7{]}) does not train the recurrence. The partition this
program registers (§4, PR-2/PR-6) is that a method earns the
in-situ-training claim only if it updates the pole-defining set
\(\{\delta_j,\ \kappa_{\mathrm{tot},j},\ \mu_{jk}\}\) on the device, and
every method in the bake-off trains the same partition. Notably, the
minimal set is \emph{gain-free} and all-thermo-optic --- heaters and
tunable couplers only --- which is what makes it drift-stable and
realistic for perturbative (SPSA-class) training on SiN; no active gain
stage is required for the headline claim.

The mapping, then, is not an analogy but an identification: a coupled
SiN ring bank \emph{is} a diagonal-complex-pole SSM with structured
coupling, its loss \emph{is} the model's damping parameter, its
realizable pole region is bounded and known, and every
recurrence-defining parameter has a physical actuator. Whether those
parameters can be \emph{trained through the physics} --- with realistic
gain saturation, ASE noise, and measurement cost --- is the question the
rest of this paper is built to answer.

\subsection{3. A pre-registered dissipative
substrate}\label{a-pre-registered-dissipative-substrate}

\subsubsection{3.1 Why one shared
substrate}\label{why-one-shared-substrate}

The bake-off's conclusions are comparisons, and comparisons inherit the
honesty of their common ground. All four training routes therefore run
on a \emph{single} simulated substrate --- one dynamical model, one
noise model, one parameter partition, one initialization and data
convention (PR-6) --- frozen and signed before any estimator existed.
Nothing in the substrate was chosen with knowledge of which method it
would favor; the freeze order is auditable in the ledger.

The model is the 2N-mode coupled-mode-theory (CMT) lattice of §2: N
rings, each carrying a clockwise/counter-clockwise doublet (backscatter
coupling is always on, §3.4), with the linear readout followed by direct
detection \(y(n) = |\sum_j c_j a_j(n)|^2\) --- the substrate's only
nonlinearity apart from gain saturation. The trainable (in-situ)
partition is fixed by the architecture freeze (PR-2): per-ring detunings
\(\delta_j\), per-ring tunable bus couplings \(\kappa_{\text{ext},j}\),
and nearest-neighbor inter-ring couplings \(\mu_{j,j+1}\) --- the
parameters that \emph{define} the recurrence, as opposed to
reservoir-style approaches that train only the readout. The digital head
(an 8-lag linear FIR on \(y\)) and the fixed affine intensity encoder
bracket the photonic core; both are identical across methods, and the
head's training cadence is part of the fairness contract.

\subsubsection{3.2 Loss, gain, and the operating
point}\label{loss-gain-and-the-operating-point}

The per-ring net decay is
\(\kappa_{\text{net},j} = \kappa_i - g_j + 2\kappa_{\text{ext},j}\):
intrinsic loss, minus Er gain, plus the loading of the through/drop
ports. Gain follows the \textbf{M1 static-saturated} class: a
per-episode operating point \(g(\bar P) = g_0/(1+\bar P/P_\text{sat})\),
differentiable in the drive statistics and in \(\kappa_{\text{ext}}\)
(the coupling changes the intracavity power the medium sees), fixed
within a rollout --- valid under the registered conditions (stationary
within-episode drive statistics, pinned train/test power statistics,
quasi-static margin
\(E_\text{sym}/E_\text{sat} \sim 10^{-6}\)--\(10^{-5}\), all
{[}EV{]}-anchored; bursty inputs void M1). Every gain-bearing cell runs
at the registered ceiling \(g_\text{rt} = 0.9\times\) the intrinsic
round-trip loss --- gain never compensates port loading --- so
\(\kappa_\text{net} = 0.1\kappa_i + 2\kappa_\text{ext}\) at the
operating point, with the passive floor (\(g=0\)) as a labelled
sensitivity row. A one-off rate-equation integrator (the salvaged M2)
validated the quasi-static reduction at the gating cells.

Amplified spontaneous emission enters as a continuous Langevin term
(convention A2),
\(\langle FF^*\rangle = 2\kappa_g n_\text{sp}\,\delta(t{-}t')\) in
photon units, discretized exactly alongside the ZOH/van-Loan propagator,
with \textbf{fresh draws on every physical pass} --- noise is never
shared between passes, methods, or the forward/reverse directions (this
single convention carries much of the fairness burden, and one of RHEL's
irreversibility invariants). The headline noise figure is \textbf{NF =
7.0 dB} (\(n_\text{sp}=2.5\)). The program's original debt register
carried ``no measured NF for the flagship Er:Si₃N₄ device'' as
verification debt \#3; the S0.3-0 substrate reconnaissance (2026-06-10)
found that premise false --- the full text reports ``a noise figure of
ca. 7 dB \ldots{} at net gain of \textgreater20 dB, limited by coupling
losses'' {[}26{]} --- and NF-A = 7.0 dB was frozen \emph{to match that
measurement} (re-verified at page level during assembly, 2026-07-12).
The residue of the debt is narrower than its original form: what exists
is a single coupling-loss-limited \emph{system} NF, not an isolated
intrinsic amplifier NF, so NF ∈ \{3, 5\} remain registered sensitivity
values.

\subsubsection{3.3 The three cells}\label{the-three-cells}

{\def\LTcaptype{none} % do not increment counter
\begin{longtable}[]{@{}
  >{\raggedright\arraybackslash}p{(\linewidth - 8\tabcolsep) * \real{0.2000}}
  >{\raggedright\arraybackslash}p{(\linewidth - 8\tabcolsep) * \real{0.2000}}
  >{\raggedright\arraybackslash}p{(\linewidth - 8\tabcolsep) * \real{0.2000}}
  >{\raggedright\arraybackslash}p{(\linewidth - 8\tabcolsep) * \real{0.2000}}
  >{\raggedright\arraybackslash}p{(\linewidth - 8\tabcolsep) * \real{0.2000}}@{}}
\toprule\noalign{}
\begin{minipage}[b]{\linewidth}\raggedright
cell
\end{minipage} & \begin{minipage}[b]{\linewidth}\raggedright
platform class
\end{minipage} & \begin{minipage}[b]{\linewidth}\raggedright
\(Q_i\)
\end{minipage} & \begin{minipage}[b]{\linewidth}\raggedright
N
\end{minipage} & \begin{minipage}[b]{\linewidth}\raggedright
role
\end{minipage} \\
\midrule\noalign{}
\endhead
\bottomrule\noalign{}
\endlastfoot
C-1 & foundry floor (P-FND) & \(2\times10^6\) & 8 & Gate-ii gating
cell \\
\textbf{C-2} & demonstrated MPW class (P-AN800) & \(6.8\times10^6\) &
\textbf{32} & \textbf{bake-off headline} \\
C-3 & aspirational (P-UHQ) & \(3\times10^7\) & 128 & labelled sweep
axis; never gates \\
\end{longtable}
}

C-2 is registered as a \emph{conservative bound} on a demonstrated MPW
result (3.3 dB/m, mean \(Q_i \approx 10.8\)M {[}20{]}); our pair
(\(5.1\) dB/m, \(Q_i = 6.8\times10^6\)) is strictly worse than the
broadest-linewidth device in that paper's full 249-resonance
distribution, whose figures are independently assessed in an invited
peer commentary {[}27{]}. The residual anchor risk is geometry transfer
(the demonstrated loss is achieved \emph{by} a wide multimode Euler-bend
racetrack; our registry ring is single-mode) --- flagged by the
commentary itself and priced by a registered ×2-loss derate row. At the
operating point, C-2 holds \textasciitilde32 samples of memory at 2 GS/s
and packs the 32-ring spectrum in-band; C-1 holds \textasciitilde9.4
(covering the task's 7-tap span at the operating point; marginal only at
its passive floor).

\subsubsection{3.4 Backscatter and the
doublet}\label{backscatter-and-the-doublet}

Surface roughness couples the CW/CCW modes at rate \(\gamma\) (11.8 MHz
at C-2, damascene-clean process class; 90 MHz at C-1), splitting each
ring resonance into a doublet --- at high \(Q_i\) the splitting is
\emph{resolved}: at C-2/θ₀ the separation clears the §2.4 criterion with
\(2\gamma/\kappa_\text{net} \approx 2.4\), and at the S0.4-0
lasing-floor calibration point (clause (a),
\(\kappa_\text{net} = 0.05\,\kappa_i\)) the ratio reaches
\(\gamma/\kappa_\text{net} = 16.6\) --- so the substrate is the full
\(2N\)-mode doublet lattice \textbf{always}, not an N-mode idealization
with a correction term. (An earlier draft quoted the floor ratio at θ₀;
the operating points are distinct and both are given here.) This choice
was consequential: the S0.4-0 calibration measured the drive build-up at
the lasing floor to be \(\times0.42\) of the hold value --- the doublet
plus chain hybridization \emph{quench} the single-pole build-up that two
independent prior estimates (×37 and ×91) had presumed, mooting both.
Modelled imperfections cannot be un-modelled by a reviewer; unmodelled
ones are the subject of §8.

\subsubsection{3.5 The feasible box, the lasing boundary, and the drive
budget}\label{the-feasible-box-the-lasing-boundary-and-the-drive-budget}

The trainable couplings live in a registered box
\(r_j = \kappa_{\text{ext},j}/\kappa_i \in
[0.1, 3]\) (K4), verified to map into the realizable pole region at
every cell. With saturating gain the box's lower edge is unsafe: below
\(r^* \approx 0.134\) the ring crosses threshold and the linear rollout
diverges. The operative bound is the S0.4-0-measured
\(\mathbf{r_\text{min} = 0.1606}\) (crossing \(+\) registered margins),
enforced as the training-time clamp; the θ₀ hold value is \(r_0 = 0.3\).
Damping, in the D-LinOSS sense, is deliberately \emph{not} a frozen
cell: it is the trainable per-ring net loss the estimators explore
through \(\kappa_\text{ext}\) over this box at fixed \(g_\text{rt}\)
(the R-ii disposition), characterized separately (§6).

Drive is normalized by an intracavity-energy budget (O2): the registered
\(\bar P_0 = 1\) mW bus drive corresponds to \(E_0 = 1.069\times10^8\)
intracavity photons at θ₀, and the digital encoder is \emph{not
permitted to buy SNR} by rescaling drive power against the registered
noise cell --- the budget is held fixed across methods and across the
in-situ/offline comparison. The S0.4-0 calibration confirmed the budget
is invariant (\(1.000000\)) under the four-tap input map below.

\subsubsection{3.6 The measured input map, and what it says about
controllability}\label{the-measured-input-map-and-what-it-says-about-controllability}

A single bus port cannot train a 32-ring chain: with drive on ring 1
only, the per-ring gradient magnitude collapses with hop distance --- an
effective participating dimension of ≈3 of 32 rings. The registered
remedy is a measured, minimal multi-point input map: S0.4-0 searched tap
sets under a pre-registered every-ring controllability gate (each ring's
gradient \(\ge 10^{-3}\) of the maximum, min over five drive seeds) and
resolved \textbf{B = taps \{3, 12, 21, 30\}} at \(K=4\) (32/32 rings
pass, worst \(1.42\times10^{-3}\); no \(K\le3\) set passes). The four
E/O drive channels this costs are charged to the systems envelope (§7)
--- multi-point drive raises exactly the conversion overhead the
envelope exists to price. Two protocol rulings made during the
measurement (gate referenced to the maximum ring; min-over-seeds
robustness) both strengthen the gate and are recorded for review. One
anchor risk stays open (vii): the calibration's floor numbers are
computed on-resonance; a detuned ring at \(r_\text{min}\) can reach
\(\kappa_\text{net} = -0.48\kappa_i\) under worst-case de-saturation,
and a δ-aware floor would sit at \(0.2547\) --- the conformance check
that would settle it was measured to be not-cheap and is carried as a
labelled risk, not silently absorbed. A pre-registered endpoint
diagnostic (PR-18; S0.11) later measured the label against the trained
solutions themselves: under the same δ-aware hypothetical, three of
eight SPSA seeds end with at least one ring in the super-threshold
corner (min \(\kappa_\text{net} = -0.06\,\kappa_i\); rings at
\(r = 0.17\)--\(0.22\) with \(|\delta| = 0.65\)--\(1.0\,\kappa_i\)),
while PAT's endpoints stay sub-threshold (\(+0.31\,\kappa_i\) at
achieved detunings) and both §6 arms sit far from the corner
(\(\geq +1.3\,\kappa_i\)). The exposure lives in the hypothetical, not
the runs (the as-built substrate is δ-independent in gain and never
lases, and the measured doublet quench of §3.4 suggests the single-pole
build-up it assumes is pessimistic) --- but the label is no longer
vacuous, and every exposed endpoint sits below the δ-aware clamp
\(r \approx 0.2547\), which is therefore the concrete mitigation Stage 1
inherits. Mid-training trajectories are not stored; the measurement
binds endpoints only.

\subsubsection{3.7 The ledger as method}\label{the-ledger-as-method}

Every consequential choice above --- the gain class and its validity
conditions, the NF axis, the cells, the box, the clamp, the input map,
the drive budget --- was written into a pre-registration ledger, in most
cases signed by the PI, \emph{before} the run that consumed it, with
measured deferrals returned to the ledger as dated addenda (the \(E_0\)
value, \(r_\text{min}\), the resolved input map, the budget \(B\), the
BPTT ceiling). Supersessions retain the superseded text. The ledger, its
review trail (independent adversarial review through 2026-07-06;
single-session self-review thereafter, disclosed), and the
frozen-before-run commit hashes ship as supplementary material. We treat
this as part of the method: a bake-off whose thresholds, budgets, and
fairness conventions are set after seeing results would not support the
claims of §5.

\subsection{4. Four routes to on-chip
gradients}\label{four-routes-to-on-chip-gradients}

\subsubsection{4.1 What counts as a physical
gradient}\label{what-counts-as-a-physical-gradient}

The fairness contract (PR-6) fixes one invariant above all: \textbf{a
training method may obtain gradient information only from simulated
device passes on the shared substrate, with fresh noise on every pass.}
Autodifferentiation through the substrate is reserved for the BPTT
reference --- the ceiling, never a contestant. All methods share, per
seed: the same initialization (detunings spread over
\([-\kappa_i, \kappa_i]\), \(r_0 = 0.3\), connected chain
\(\mu_c = 0.3\kappa_i\)), the same data stream (a function of seed and
iteration only), the same head cadence, the same clamp to the operative
box, and the same pre-registered hyperparameters. Cost is counted in
\textbf{physical device passes, any direction} (PR-7): one pass = one
sequence through the substrate. Digital compute (a twin's
forward/backward, the head update) lives on a side-ledger that is always
co-reported and never folded into the rank --- a method that saves
device passes by spending digital FLOPs has a real but \emph{different}
advantage than a model-free one, and merging the ledgers would hide
exactly the distinction the comparison exists to draw.

\subsubsection{4.2 SPSA --- model-free, two
passes}\label{spsa-model-free-two-passes}

Simultaneous-perturbation stochastic approximation {[}1{]} perturbs the
entire in-situ partition by \(\pm c\Delta\) (a random sign vector),
measures the scalar loss twice, and forms a descent direction from the
difference: \textbf{2 device passes per update}, no model, no twin, no
added hardware beyond the plant's own actuators and its single readout
(the simplest row of the hardware ledger, supplementary note N2).
Perturbations at the clamp boundary are one-sided. SPSA is
chip-demonstrated {[}2{]} and inherits the crosstalk-robustness observed
in our prior thermo-optic work {[}28{]}; its known weakness --- gradient
variance growing with parameter count --- is precisely what the
sample-efficiency metric prices.

\subsubsection{4.3 PAT --- physical forward, twin
backward}\label{pat-physical-forward-twin-backward}

Physics-aware training {[}3{]} evaluates the loss on the \emph{measured}
physical output and takes the parameter gradient through a
differentiable digital twin at the commanded parameters: \textbf{1
device pass per update}, plus a twin forward/backward on the digital
ledger. PAT's value proposition is exactly this exchange, and its
honesty hinges on the twin being \emph{imperfect in a registered way}.
The twin-mismatch protocol (PR-5) freezes three families: \textbf{M-par}
--- 5\%-class parametric calibration errors on the constants and
actuation maps (\(\kappa_i{+}5\%\), \(\gamma{+}5\%\), actuation
×1.05/×0.95, detuning offset \(0.05\kappa_i\), a loose gain pair
×1.10/×0.75 reflecting debt \#3); \textbf{M-struct} --- a structural
omission (the twin drops the gain's dependence on the trained coupling,
the one channel a fixed-gain model cannot see); and \textbf{M-noise}
(always on) --- the twin is noiseless. The headline PAT is the
\emph{composed} mismatch; the decomposition is reported separately
(§5.6). A perfect-twin gate verifies the implementation: with mismatch
off, PAT's gradient equals BPTT's to machine precision.

\subsubsection{4.4 Recurrent in-situ adjoint --- a physical reverse
pass, by
hypothesis}\label{recurrent-in-situ-adjoint-a-physical-reverse-pass-by-hypothesis}

The adjoint route extends feedforward in-situ backpropagation {[}4{]} to
the time-domain cavity setting: the error field is physically propagated
\emph{backward} through the same dissipative substrate, and the gradient
is read from forward/adjoint interference. No recurrent photonic
demonstration of this pass exists (verification debt \#4); the
simulation charges it \textbf{as if realizable} --- 1 forward + 1
adjoint = \textbf{2 device passes per update}, zero digital --- and
quarantines the realizability question in the hardware ledger
(circulators, phase-coherent injection, separability of the
counter-propagating field from the backscatter doublet). The simulated
adjoint pass is honest about two physical limits: it carries
\textbf{fresh noise} (the reverse pass is its own noisy traversal), and
the saturating gain is \textbf{frozen at its operating-point value} in
the backward linearization --- a counter-propagating field experiences
the medium's saturation state but cannot realize the
\(\partial g/\partial\kappa_\text{ext}\) self-consistency channel. Both
simplifications flatter the adjoint (a third, an additive noise term on
the adjoint field itself, awaits an error-launch power convention that
is unresolved hardware design), so the bake-off's adjoint arm is an
\textbf{optimistic bound} and is labelled as such wherever it appears.
Floor checks: on the fixed-gain plane the adjoint gradient equals BPTT
to machine precision; in saturating mode the frozen-gain approximation
costs 0.6\% of gradient direction at C-1 but \textbf{7.5\% at C-2} ---
the one method-relevant quantity we measured to \emph{grow} with cell
size.

\subsubsection{4.5 RHEL --- Hamiltonian echoes on a substrate that
forgets}\label{rhel-hamiltonian-echoes-on-a-substrate-that-forgets}

Recurrent Hamiltonian echo learning {[}29, 30{]} trains by
time-reversal: evolve forward; apply a single conjugation to the state
snapshot (for optical fields, phase conjugation); evolve again through
the \emph{same} physics with the input replayed time-reversed and a
small error nudge injected continuously; the gradient is the symmetric
finite difference of \(\nabla_\theta H\) between a \(+\varepsilon\) and
a \(-\varepsilon\) echo. Three consequences of taking the published
algorithm seriously on a dissipative substrate, each registered before
the runs: (i) the count is \textbf{3 passes in the Hamiltonian limit but
4 operationally} (2 forward + 2 echo) --- the echo does not return the
state re-usable, and state cloning is unphysical (PR-7.1); (ii) the
conjugation fires \textbf{twice per update}, paying its penalty chain
independently each time; (iii) the update rule reads only the
\emph{coherent} generator, so the dissipative channel of
\(\kappa_\text{ext}\) is structurally invisible to it.

The echo's physical primitive is modelled concretely (PR-11), not as an
idealized operator: a \(\chi^{(3)}\) four-wave-mixing conjugation stage
with the full penalty chain --- extraction through the ring ports
(\(\eta_\text{ex} = 2\kappa_\text{ext}/\kappa_\text{net} = 0.86\) at
θ₀), single-pass spiral conversion
(\((\gamma_\text{nl} P_p L)^2 \approx -16.7\) dB at 0.3 W pump, 0.5 m;
\(\gamma_\text{nl} \approx 0.97\,\text{W}^{-1}\text{m}^{-1}\) for
tight-confinement SiN {[}31{]}), and timing decay --- totalling
\textbf{−22.4 dB per conjugation} (Fig. S3), plus the phase-insensitive
parametric quantum floor. A page-level check at assembly found published
\emph{ultra-low-loss-geometry} demonstrations at
\(\gamma \approx 0.29\)--\(0.51\,\text{W}^{-1}
\text{m}^{-1}\) (with CW power handling demonstrated to 7 W) {[}31{]};
our 0.97 assumes a tighter-confinement spiral than those demos, so the
frozen chain is, if anything, \emph{optimistic} for RHEL --- at the
measured ULL values the conjugation penalty deepens by a further
\textasciitilde6 dB. The direction only strengthens §5.4's conclusion,
which the idealized-conjugator control shows does not hinge on the chain
at all. The quantum floor itself was measured negligible at the
\(10^5\)-photon state scale (as registered-to-measure), with
pump-transfer excess included. Because the ring fields overlap
spectrally, conjugating N rings needs N pumped arms: \textbf{≈9.6 W of
on-chip pump at the headline cell}, charged to the envelope. Off-chip
conjugation was priced and excluded (the state decays in transit:
amplitude survival 0.12--0.54 at 10 ns for C-1/C-2; there is no storage
primitive to wait out a conjugator). Three irreversibility invariants
bind the implementation and are test-enforced: independent forward/echo
noise streams (no common-RNG reversal); no loss-sign flip (the echo
traverses the same dissipative lattice); gain injects fresh ASE in the
echo too --- and the echo of a noisy forward must \emph{not} recover the
noiseless state. The floor check completes the picture: as dissipation
is removed, the implemented RHEL gradient converges to the exact
reference (cosine \(\to 1.0000\)), so whatever §5.4 finds is the
physics, not the code.

\subsubsection{4.6 The guardrail}\label{the-guardrail}

The four routes are \emph{parallel in simulation, singular in hardware}:
the Stage-1 chip is committed to the two chip-demonstrated workhorses
(PAT, SPSA) regardless of the bake-off's ranking, and an exact-gradient
route can earn a \emph{later} hardware slot only by clearly beating both
on a pre-registered outcome criterion (PR-9; ``exactness'' is struck
from the promotion menu). §5 reports how this resolved.

\subsection{5. The bake-off: which routes train the recurrence, and at
what
cost}\label{the-bake-off-which-routes-train-the-recurrence-and-at-what-cost}

\subsubsection{5.1 Pre-registration and the reference
ceiling}\label{pre-registration-and-the-reference-ceiling}

Every threshold in this section was fixed in the ledger \emph{before}
the run it judges, in the order the runs consumed them: the fairness
contract, cost metric, and twin-mismatch families (PR-6/7/5) at the
start of S0.4; the target rule, statistical plan, and gate semantics
(PR-3/8/9) at S0.4-close; and the two measured numbers a rule cannot
supply --- the budget and the ceiling --- as dated addenda committed
before any contestant ran. We state this not as ceremony but because the
central claim is a \emph{threshold-crossing} claim, and a threshold
chosen after seeing the curve is worthless.

The reference is backpropagation-through-time on the substrate itself
(BPTT), which is not a physical training method --- it reads gradients
the device cannot expose --- and is the strongest gradient access the
substrate model admits. It is a \emph{budget-scoped reference}, not a
capacity ceiling: what BPTT reaches within a stated update budget is
what training can be held to at that budget --- and the §17.8 diagnostic
below measured BPTT still improving past the bake-off budget, so no
asymptotic-capacity claim is made, and none is needed. On the headline
cell (C-2: 32 rings, \(Q_i = 6.8\times10^6\); 4-PAM channel equalization
at 28 dB, §3), BPTT drives the symbol-error rate to a median of
\(5.2\times10^{-4}\) across eight seeds (seven of eight at
\(5\times10^{-4}\) --- two errors in the 3840-symbol evaluation set, the
quantization floor). Because that floor is too coarse for comparisons
\emph{between} near-ceiling arms --- a two-symbol band can manufacture
or erase a ``statistically real'' difference --- every such comparison
in this section (§5.5, §5.6) was re-scored under a pre-registered fine
protocol (\textbf{eval-F}, PR-17: the same reserved held-out streams
extended to 99,840 scored symbols, per-seed resolution
\(1.0\times10^{-5}\)), with the frozen coarse protocol remaining the
protocol of record for every gate, target, and ranking verdict, and both
reported wherever they differ. At eval-F the BPTT reference itself reads
\(9.8\times10^{-4}\) (the coarse \(5.2\times10^{-4}\) was a
lucky-two-errors reading --- an illustration of exactly the floor
hazard). One scope note, so no reader has to discover it: this reference
is \emph{protocol-local}. It is measured under the bake-off contract ---
12,000 updates at the registered cadence --- and it bounds the §5.3
ranking, which runs under that same contract. The follow-up experiments
of §5.5 run their own frozen specification (PR-5 §E) with a
\(2.6\times\) larger update budget (31,600), and at eval-F their trained
arms land slightly \emph{below} the 12,000-update number
(\(8.4\)--\(9.0\times10^{-4}\) vs \(9.8\times10^{-4}\); the paired
inversion is within seed noise, CI \([-0.1, +3.1]\times10^{-4}\)
including zero). A reference matched to the follow-up budget --- BPTT at
31,600 updates, registered as a review-added diagnostic with its
expected direction stated in advance (PR-17 §17.8) --- lands at
\(8.1\times10^{-4}\) (8 seeds, per-seed \(7.3\)--\(8.9\times10^{-4}\)),
and it, not the bake-off number, is the reference line drawn in Fig. F8.
We state the margin rather than leave it to be measured: the restoration
of the expected ordering is \emph{median-level and thin} --- the matched
reference sits \(0.3\)--\(0.9\times10^{-4}\) below the calibration-sweep
arms and one per-seed resolution unit below the drift arm's
time-integrated \(8.2\times10^{-4}\), with its per-seed spread
overlapping the arms' --- a scope statement (with matched budget, no
physical arm sits meaningfully below exact gradients), not a separation.
An arm sitting below the 12,000-update number is a budget effect, not a
physical estimator beating exact gradients. One implementation erratum
in the eval-F tooling (an evaluation-normalization inconsistency, caught
the same day by its chance-level signature, registered, fixed under a
machine-precision gate test, and re-run) is documented in the ledger
(PR-17 §17.7). \textbf{The frozen substrate has headroom for the task;
the open question is purely which physical training routes reach it, and
at what cost.} The pre-registered target follows mechanically: a method
\emph{reaches target} if its held-out SER falls to
\(\text{SER}_\text{target} = 1.25\times\text{ceiling} + 0.005 = 5.65\times
10^{-3}\) at any evaluation point within the device-pass budget (the
additive guard dominates at a floor-level ceiling, by design --- PR-3
§B; had the rule consumed the eval-F ceiling instead, the target would
be \(6.23\times10^{-3}\) and no reach-target verdict in this paper
changes). The budget \(B = 252{,}800\) device passes is twice the BPTT
convergence point measured in a seed-7 sizing pilot excluded from the
eight scored seeds --- where ``convergence'' means the pilot's
registered flatness rule, a budget-local criterion: the §17.8 matched
run kept improving at the fine floor past that point, which the flatness
rule (defined on the coarse trace) could not resolve.

A note this cell settles for free: the ceiling is \emph{identical} under
fixed-gain and saturating-gain substrate models (\(\Delta_{M3}=0\)). The
gain-model class --- the largest modelling uncertainty in the substrate
(§3) --- cannot move the achievable accuracy, so it cannot flip any
ranking below. The pre-registered M3 sensitivity trigger is therefore
un-triggerable at this cell.

\subsubsection{5.2 Gate ii: the recurrence trains on the
device}\label{gate-ii-the-recurrence-trains-on-the-device}

Two decompositions of the Stage-0 gate were pre-registered (PR-9).
\textbf{Capacity (ii-a):} the ceiling must clear a task-utility floor
set at half the readout-only error --- the reservoir baseline that
freezes the recurrence and trains only the digital head. The baseline
stalls at \(2.2\times10^{-2}\) (scored at the coarse protocol; at
\textasciitilde85 errors per 3,840-symbol evaluation it sits far from
the quantization floor, so eval-F cannot move it materially), putting
the floor at \(1.1\times10^{-2}\); the reference clears it by
\(21\times\) at the coarse protocol of record (\(5.2\times10^{-4}\)) and
by \(11\times\) at eval-F (\(9.8\times10^{-4}\)). \textbf{Trainability
(ii-b):} at least one of the two hardware-committed routes ---
physics-aware training (PAT) or SPSA --- must reach target on at least
five of eight seeds.

\textbf{Both reach target on all eight.} This is the load-bearing result
of the program: the parameters that \emph{define} the recurrence --- the
per-ring detunings, the tunable ring--bus couplings, and the inter-ring
couplings \(\{\delta_j, \kappa_{\text{ext},j}, \mu_{jk}\}\) --- are
trained on the (simulated) physical substrate, through
physical-operation-only gradient methods with fresh injected noise on
every pass, to within the pre-registered margin of the exact-gradient
ceiling. To our knowledge this is the first demonstration --- \textbf{in
simulation, on a pre-registered realistic substrate model} --- of a
continuous-time dissipative-resonator recurrence whose poles and
couplings are updated by device-protocol gradient-based or
gradient-estimating training (the white-space claim, §1; {[}32{]}). The
on-chip counterpart does not exist yet; it is what Stage 1 is designed
to earn (§9), and every use of ``demonstration'' in this paper carries
this qualifier.

\subsubsection{5.3 The ranking: sample-efficiency at matched device-pass
cost}\label{the-ranking-sample-efficiency-at-matched-device-pass-cost}

The primary metric is sample-efficiency: device passes to reach target,
one physical pass being one sequence through the substrate in any
direction (PR-7). Ranked lexicographically by success fraction then
median passes (PR-8), with the digital-compute side-ledger co-reported
but never folded into the rank (converting both ledgers to joules
\emph{inverts} this ranking --- §7.3; the two-ledger principle exists
precisely so that neither metric is presented as the truth):

{\def\LTcaptype{none} % do not increment counter
\begin{longtable}[]{@{}
  >{\raggedright\arraybackslash}p{(\linewidth - 8\tabcolsep) * \real{0.2000}}
  >{\raggedright\arraybackslash}p{(\linewidth - 8\tabcolsep) * \real{0.2000}}
  >{\raggedright\arraybackslash}p{(\linewidth - 8\tabcolsep) * \real{0.2000}}
  >{\raggedright\arraybackslash}p{(\linewidth - 8\tabcolsep) * \real{0.2000}}
  >{\raggedright\arraybackslash}p{(\linewidth - 8\tabcolsep) * \real{0.2000}}@{}}
\toprule\noalign{}
\begin{minipage}[b]{\linewidth}\raggedright
route
\end{minipage} & \begin{minipage}[b]{\linewidth}\raggedright
success
\end{minipage} & \begin{minipage}[b]{\linewidth}\raggedright
median device passes → target
\end{minipage} & \begin{minipage}[b]{\linewidth}\raggedright
final SER (median)
\end{minipage} & \begin{minipage}[b]{\linewidth}\raggedright
digital ledger (at budget \(B\))
\end{minipage} \\
\midrule\noalign{}
\endhead
\bottomrule\noalign{}
\endlastfoot
\textbf{PAT} (twin-backward) & 8/8 & \textbf{38,400} &
\(8\times10^{-4}\) & 505,600 \\
\textbf{adjoint}† (physical reverse pass) & 8/8 & 73,600 &
\(5\times10^{-4}\) & 0 \\
\textbf{SPSA} (model-free) & 8/8 & 176,000 & \(1.3\times10^{-3}\) & 0 \\
\textbf{RHEL}‡ (Hamiltonian echo) & 0/8 & censored at \(B\) &
\(1.4\times10^{-1}\) & 0 \\
\end{longtable}
}

† \emph{Charged as-if-realizable: no recurrent physical reverse pass has
been demonstrated on any platform (debt \#4, §8.2); this row is an
optimistic bound on a hypothetical implementation and competes with
chip-proven routes only in that idealized sense (§4).} ‡ \emph{Reported
as a measured feasibility bound on echo learning in dissipative
substrates, not as a competitive entry --- see §5.4.}

All three ordered pairs among the passing routes separate with
paired-by-seed bootstrap 95\% confidence intervals excluding zero
(PAT−SPSA \(=-137{,}600\) passes, CI \([-155{,}200,-123{,}200]\);
adjoint−PAT \(=+35{,}200\), CI \([+35{,}200,+36{,}800]\); adjoint−SPSA
\(=-102{,}400\), CI \([-120{,}000,-88{,}000]\)). The adjoint−PAT
interval is degenerate --- its point estimate sits on its own lower
bound, one 1,600-pass step wide --- and we say why rather than let it
read as precision: passes-to-target lives on a 1,600-pass evaluation
grid (PR-3's no-interpolation rule), the eight paired differences all
fall within two grid steps (\(+33{,}600\) to \(+36{,}800\), positive on
8/8 seeds), and a bootstrap over values that concentrated collapses onto
the grid. It should be read as a sign-consistent separation bounded by
the grid resolution, not as a distributional interval. The three routes
trade the same axes the theory predicts they should. \textbf{PAT} is
cheapest on the physical device but pays for it digitally --- two twin
passes (forward + backward) for every physical pass, a side-ledger of
505,600 digital passes over the full budget (the co-reported column
above; the device-pass column is to-target, so the two columns are
deliberately not a ratio) --- and carries the full burden of
characterizing a differentiable twin (§4; hardware ledger, supplementary
note N2). \textbf{The adjoint} matches the exact-gradient ceiling in
accuracy at zero digital cost and \(1.9\times\) PAT's device passes ---
though its count charges one physical reverse pass \emph{as if}
realizable, which no recurrent photonic system has yet demonstrated (a
caveat we quarantine, §4). \textbf{SPSA} costs \(4.6\times\) PAT's
device passes but needs no model, no twin, and no added hardware --- the
simplicity anchor of the Stage-1 plan.

No route earns promotion toward a later hardware slot (PR-9): the
adjoint beats SPSA but loses to PAT on device passes, and clearing the
bar requires clearly beating \emph{both} workhorses. The hardware
roadmap therefore stays on PAT and SPSA --- the outcome the guardrail
was built to protect, now settled by data rather than assertion.

\subsubsection{5.4 RHEL: a feasibility bound on echo learning in
dissipative
substrates}\label{rhel-a-feasibility-bound-on-echo-learning-in-dissipative-substrates}

The fourth route, recurrent Hamiltonian echo learning (RHEL), is best
read not as a contestant but as a \emph{measured feasibility bound} ---
and we say plainly that its headline outcome was foreseeable in
direction, if not in magnitude, before the run: the theorem behind the
echo assumes a non-dissipative system, and the headline operating point
sits at \(\kappa_\text{net} T\, dt \approx 8\) (the C-1 control cell at
\(\approx 27\)) --- one to two orders beyond the \(\lesssim 0.1\) regime
where our own recovery curve shows the update aligning with the true
gradient (Fig. S5). (An earlier draft quoted the C-1 ratio at the
headline cell; both are given here.) What the bake-off adds is the
\emph{quantified boundary} --- where echo learning breaks on a
dissipative substrate, by how much, and through which mechanism ---
under the same fairness contract as the routes that pass; that, not a
horse race it could not win, is the result we consider citable. RHEL is
also the route whose physical primitive SiN is least suited to supply.
Rather than an idealized conjugation operator, we model the echo as a
concrete \(\chi^{(3)}\) four-wave-mixing phase-conjugation stage with
its measured penalty chain --- extraction, single-pass spiral conversion
at 0.3 W pump, timing decay --- totalling \(-22.4\) dB per conjugation,
plus the phase-insensitive parametric noise floor (§4, PR-11; {[}31{]}).

RHEL does not reach target on any seed; its final SER of \(0.14\) is
\emph{worse than the readout-only baseline} (a \(+0.118\) readout
differential). Two controls locate the cause. A floor check confirms the
estimator is correct: as the substrate is made progressively less
dissipative, the RHEL gradient converges to the exact reference
(direction cosine \(\to 1.0000\); Fig. S5) --- the non-dissipative limit
RHEL's theorem assumes. And an idealized-conjugator control at the
smaller C-1 cell --- a \emph{perfect} echo, no conjugation loss or noise
--- formally reaches the C-1 target
(\(5.7\times10^{-3} \le 7.3\times10^{-3}\) coarse; \(7.2\times10^{-3}\)
at eval-F, §5.6) but converges to the head-only level: the recorded
S0.4c diagnosis is that even a perfect echo's \emph{recurrence}
contribution is \(\approx 0\) in this regime --- the digital head does
the passing. The conjugation chain is therefore second-order, and the
failure at the headline cell is dissipative-echo bias itself: the
irreducible mismatch between an echo that assumes time-reversal and a
substrate whose 256-sample sequence spans \(\approx 8\) amplitude memory
lifetimes (memory \(\approx 32\) samples at C-2; the honest echo's
updates are then not merely useless but harmful --- the \(+0.118\)
readout differential above). This is the honest instantiation of the
platform argument (§1, §5.4 of the proposal): silicon nitride's low loss
improves RHEL's \emph{noise} budget, but the \emph{dissipation} the
recurrence itself requires is fatal to the echo at the operating point.
RHEL-on-SiN stays a simulation result.

\subsubsection{5.5 The comparison the fair design was built to
expose}\label{the-comparison-the-fair-design-was-built-to-expose}

One baseline result is more consequential for the program than the
ranking. The offline-train-then-deploy route --- train the full
parameter set digitally on a designer's model, then deploy through
actuation maps, recalibrating only the digital head on-device --- was
given the \emph{same} 5\%-class calibration errors as PAT's twin (the
mismatch families are drawn from one frozen set, so the comparison
cannot be rigged by giving in-situ training a secretly-wronger
competitor; PR-5). At that mismatch level it reaches
\(1.0\times10^{-3}\) against in-situ PAT's \(8\times10^{-4}\) (coarse
protocol of record --- a one-to-two-symbol gap the coarse floor scores
as formally real; at eval-F, where the floor can actually resolve it,
the two arms show no resolved difference, CI including zero --- the F8a
sweep below).

We state the consequence plainly, because the fair-comparison design
exists precisely to force it: \textbf{at 5\% calibration accuracy on
this task, we resolve no performance benefit for in-situ training over
calibrate-then-deploy.} This does not establish equivalence or rule out
smaller benefits. The result established here is simulated trainability
of the recurrence-defining parameters; a physical demonstration remains
future work. The question this raises --- \emph{under what conditions
does the advantage appear?} --- we then answered with two pre-registered
follow-up experiments rather than leaving it open (PR-5 §E, PR-16; both
frozen before the runs).

\textbf{Corrected calibration sweep: no resolved advantage through the
registered 30\%-class point (S0.13; Fig. F8a).} The original sweep
failed to scale three of PAT's five registered mismatch terms in command
binding. Its higher-mismatch comparisons were withdrawn, and the code
was corrected. A bounded correction protocol was committed before the
rerun (\texttt{bda989f}): 32 fresh PAT units at m=\{2,3,4,6\}, with all
eight original seeds, unchanged training budgets and PR-17 evaluation.
The 40 unchanged offline units and eight PAT m=1 units were reused. An
additional m=1 seed-11 run reproduced every coarse evaluation, final
coarse/fine SER, and pass ledger exactly. This anchor is a
reproducibility check, not a ninth replicate. The correction is
post-result bug repair, not a newly blinded experiment; the thresholds
and bootstrap rule were not changed. Results are recorded at
\texttt{174558c}.

At eval-F (99,840 symbols), all five paired-by-seed difference intervals
include zero. Offline/PAT median ratios range from 0.994 to 1.053;
neither arm's median fails the registered accuracy target. No mismatch
crossover clears the frozen rule (ratio ≥2 and difference-CI lower bound
\textgreater0). These data resolve no difference at the tested levels;
they do not establish statistical equivalence.

{\def\LTcaptype{none} % do not increment counter
\begin{longtable}[]{@{}
  >{\raggedright\arraybackslash}p{(\linewidth - 8\tabcolsep) * \real{0.1667}}
  >{\raggedleft\arraybackslash}p{(\linewidth - 8\tabcolsep) * \real{0.2222}}
  >{\raggedleft\arraybackslash}p{(\linewidth - 8\tabcolsep) * \real{0.2222}}
  >{\raggedleft\arraybackslash}p{(\linewidth - 8\tabcolsep) * \real{0.2222}}
  >{\raggedright\arraybackslash}p{(\linewidth - 8\tabcolsep) * \real{0.1667}}@{}}
\toprule\noalign{}
\begin{minipage}[b]{\linewidth}\raggedright
mismatch class
\end{minipage} & \begin{minipage}[b]{\linewidth}\raggedleft
PAT median SER ×10⁻³
\end{minipage} & \begin{minipage}[b]{\linewidth}\raggedleft
offline median SER ×10⁻³
\end{minipage} & \begin{minipage}[b]{\linewidth}\raggedleft
offline/PAT
\end{minipage} & \begin{minipage}[b]{\linewidth}\raggedright
paired 95\% CI of (offline−PAT) ×10⁻³
\end{minipage} \\
\midrule\noalign{}
\endhead
\bottomrule\noalign{}
\endlastfoot
5\% & 0.851 & 0.896 & 1.053 & {[}-0.050, 0.100{]} \\
10\% & 0.861 & 0.891 & 1.035 & {[}-0.060, 0.120{]} \\
15\% & 0.861 & 0.866 & 1.006 & {[}-0.045, 0.080{]} \\
20\% & 0.861 & 0.881 & 1.023 & {[}-0.045, 0.100{]} \\
30\% & 0.881 & 0.876 & 0.994 & {[}-0.090, 0.070{]} \\
\end{longtable}
}

The coarse protocol is co-reported: ratios are 1.33--1.50, with positive
difference intervals at 5--20\% and an interval touching zero at 30\%;
none reaches the 2× advantage threshold. Coarse and fine metrics in each
corrected unit use the same trained model and normalization. The old
higher-mismatch outputs remain invalid historical records, not
replications of the corrected runs. Supplementary N8 and
\texttt{results/s0\_13/analysis.json} retain the correction and
input-hash trail. The m=1 bake-off and separate drift experiment are
unaffected by the binding defect.

\textbf{Drift is the axis --- specifically the part a re-lock cannot
catch} (Fig. F8b,c). We then let the substrate \emph{drift}: a random
walk on the ring detunings calibrated to a measured free-running
silicon-nitride resonance drift (\(\approx 341\) MHz over 24 h
\(\approx 24\,\kappa_i\) at C-2 {[}33{]}), deployed after convergence,
with each arm allowed its on-device response --- offline recalibrates
the head and re-locks the laser (a single global detuning re-centering);
in-situ retrains the recurrence. (The re-lock is not a strawman: it is
the strongest response available without per-ring observability --- any
per-ring re-trim requires per-ring on-device measurement and actuation
feedback, which is the in-situ stack by another name; the registered
arms are the two coherent extremes.) The pre-registered contrast holds
cleanly, and at the eval-F protocol of record the verdict is now formal.
Under \textbf{common-mode} drift (whole-chip thermal wander) the laser
re-lock absorbs it and offline keeps pace (\(0.98\) vs
\(0.73\times10^{-3}\) time-integrated, ratio \(1.34\), CI including zero
--- no advantage). Under \textbf{independent} per-ring drift the re-lock
\emph{cannot} fix the scrambled pole scatter, and in-situ retraining
pulls ahead: \(0.82\times10^{-3}\) versus the re-locking offline's
\(1.98\times10^{-3}\) time-integrated. The frozen decision rule ---
PR-16, ratified 2026-07-24 \emph{before any drift run existed}, and
re-applied verbatim at eval-F by PR-17 §17.3 --- declares an advantage
iff the offline median is \(\ge 2\times\) the in-situ median
\textbf{and} the paired-by-seed bootstrap 95\% CI of the difference
excludes zero. \textbf{Both conditions hold (ratio \(2.42\); difference
CI \([+0.71, +2.84]\times10^{-3}\)): this is the one comparison in the
paper where in-situ training formally beats the strong offline
baseline.} Because that rule joins a point-estimate threshold to a
difference-CI, we also report what it does not itself guarantee: the
\emph{ratio's} own bootstrap CI is \([1.7, 4.6]\), including values
below 2 (17\% of resamples) --- the declared advantage is a
threshold-crossing under a rule frozen in advance, not a 95\%-confidence
claim that the true ratio exceeds 2.

A verdict that flips from failed (coarse, \(1.84\times\)) to passed
(fine, \(2.42\times\)) under a protocol change is the classic post-hoc
pattern, so we dismantle it factor by factor. First, the movement is
\emph{not} confounded with a fresh noise draw. The eval-F drift numbers
come from re-executing the frozen S0.9b unit specifications, and the
re-executed trajectories are \textbf{bit-identical} to the originals:
the dynamics are deterministic given the registered seed, evaluation
draws only from reserved streams that never touch the training
randomness (the coarse evaluation set is the first 2 of the fine
protocol's 52 batches --- a strict subsample, so the coarse reading was
underpowered, not contradicted), and the evaluation-independent
device-state fingerprint matches the stored originals exactly on all 64
drift units --- as did the coarse final SER on all 80 historical
calibration-sweep units before the S0.13 binding correction (ledger
§17.8). The coarse→fine change therefore carries exactly one factor, the
evaluation floor: the \emph{same physical trajectories} score
\(1.84\times\) at 3,840 symbols and \(2.42\times\) at 99,840. Second,
the ordering, with dates and commits rather than assertion: the decision
rule froze before any drift datum existed (PR-16, 2026-07-22/24), and
the coarse estimate --- below the bar --- was on record when the fine
protocol was registered (PR-17, commit \texttt{241204a}, 2026-07-27,
\emph{before any eval-F measurement}; the same commit froze the rule
that coarse and fine verdicts are both reported wherever they differ).
Third, at the unaffected m=1 calibration point, the finer evaluation
erased a coarse-floor difference in in-situ's favor. The higher-mismatch
correction above is a separate intervention and is not used to claim
bit-identity with the invalid historical implementation. The drift
comparison crossed the registered threshold at the finer resolution. The
gap \emph{grows with accumulated drift} (\textasciitilde3--4× at the
largest drift step), exactly as the mechanism predicts. Two protocol
notes for honesty: the in-situ SPSA arm is plotted for context only ---
its per-step re-convergence transient (and, within eval-F, a registered
per-step scale-re-measure convention that penalizes an arm whose
couplings move during the step) inflate its early trajectory, and no
registered verdict involves it; and the deploy-time
\(\mathrm{SER}(t{=}0)\) diagnostic of the in-situ arms shares that
convention and is not used in any comparison.

The corrected comparisons show no resolved calibration difference at the
five tested 5--30\%-class levels and no common-mode drift advantage,
alongside a declared \(2.42\times\) advantage under independent per-ring
drift. The relevance of the drift advantage to hardware depends on the
correlation of actual on-chip drift, which has not been measured for
this architecture.

\subsubsection{5.6 What the diagnostics
add}\label{what-the-diagnostics-add}

Two mechanism rows, at three seeds each on C-1, support the mismatch
narrative without inflating it (Fig. S2; eval-F). Decomposing PAT's twin
mismatch --- perfect twin, parametric-error twin, structural-omission
twin (dropping the gain self-consistency channel) --- the perfect and
M-struct twins are indistinguishable: identical medians
(\(1.11\times10^{-3}\)), per-seed values agreeing to within
\(2\times10^{-5}\) (two symbols even at eval-F). The fine floor resolves
a small M-par excess (\(1.31\times10^{-3}\) median, +18\% relative; the
paired excess is positive on all three seeds, \(+0.2\) to
\(+2.0\times10^{-4}\), though at \(n=3\) that is a consistent sign, not
a confidence interval --- invisible at the coarse floor, where all three
arms had read as one number): PAT absorbs the structural omission
completely and the parametric family almost completely at this cell. We
flag explicitly that this does \textbf{not} extrapolate to C-2, where
the dropped gain channel was measured to carry \textasciitilde8\% of the
gradient direction (§4, the adjoint cosine dropping from 0.994 to 0.925
with cell size); the C-2 mismatch sensitivity is a Stage-0.5-full
measurement, not an inference from C-1. The idealized-RHEL row is the
control cited in §5.4 --- at eval-F it still clears the C-1 target,
barely (\(7.2\times10^{-3} \le 7.3\times10^{-3}\)), which we note
because a two-symbol coarse floor could not have resolved how thin that
margin is. Neither row is a headline; both are the pre-registered
controls that let the headlines mean what they say.

\subsubsection{5.7 Controllability: what ``N = 32'' actually
means}\label{controllability-what-n-32-actually-means}

The bake-off's cell label understates a constraint that any hardware
implementation inherits, so we report it as a first-class result (Fig.
F3). Under a single input tap, the per-ring gradient magnitude collapses
geometrically with distance from the drive --- by ring 32 it sits some
twenty-five orders of magnitude below the maximum --- and the
\emph{participation profile} (settled per-ring amplitude relative to the
maximum) counts only \(\{1, 3, 5\}\) of 32 rings above
\(\{10^{-1}, 10^{-2}, 10^{-3}\}\). A nominally 32-ring lattice driven at
one port is, effectively, a three-ring computer with 29 passengers --- a
controllability property of the chain physics, not of any training
method. (The §5.2 readout-only baseline is a \emph{separate} falsifier,
not this one's consequence: that baseline runs under the resolved
four-tap map --- full amplitude participation --- and still stalls at
\(2.2\times10^{-2}\). What it lacks is trained poles, not drive
coverage; an earlier draft wrongly attributed its stall to single-tap
starvation.)

The pre-registered remedy is a measured, minimal input map: the smallest
tap set (capped at \(K = 4\)) under which \emph{every} ring's gradient
clears \(10^{-3}\) of the maximum. The resolved map, taps
\(\{3, 12, 21, 30\}\), clears the gate for all 32 rings with worst ratio
\(1.42\times10^{-3}\) --- taken as a \emph{minimum over five drive
realizations}, because single-seed margins at the gate boundary flicker
by a factor of \textasciitilde20. No three-tap set clears (best: 24 of
32), and the registered starting guess \(\{1, 9, 17, 25\}\) was not the
winner (28 of 32 --- its worst ring sat seven hops from a tap). Under
the resolved map the participation counts rise to \(\{4, 26, 32\}\), and
every use of ``\(N = 32\)'' in this paper carries that measured profile
rather than the nominal dimension. Two protocol rulings made during
resolution are on record: the gate references the
\emph{maximum-gradient} ring (the registered ring-1 reference is
gameable when ring 1 is untapped), and robustness binds on the
min-over-seeds. Both strengthen the gate; both were adopted before the
finalists were evaluated. The price of controllability is charged
honestly where it lands: four drive E/O channels instead of one, priced
in the systems envelope (§7) --- trainability of the deep lattice is
bought with exactly the conversion overhead the advantage question
(§7.4) must then carry.

Resolved at θ₀, the map's guarantee is \emph{a priori} for the
initialization only --- so a pre-registered diagnostic (PR-18; S0.11)
reproduced the winning routes' converged solutions bit-identically (and
both §6 arms'; 32/32 seed-runs eval-trace-exact) and re-ran this
section's gate \emph{at the solutions}. The demonstration routes
maintain it: PAT clears \textbf{32/32 on every seed} with worst ratio
\(1.7\times10^{-3}\) --- the θ₀-class margin --- and SPSA holds a median
of 30.5/32 (range 29--32). The mechanism is the one §6 identifies: both
winning routes, on every seed, converge to the \emph{same} actuator
structure --- the four driven rings damped to \(r \approx 1.3\) (median
1.29 PAT, 1.32 SPSA) while the undriven rings sit \emph{at} their
initialization (median \(r \approx 0.30\) against \(r_0 = 0.3\),
within-band sd 0.05--0.07; the whole-profile
\(\mathrm{sd}(r) \approx 0.33\) reported in S0.11 is carried almost
entirely by the four-tap excursion) --- heterogeneity that buys damping
where the input lands without severing gradient transport. The §6 boxed
arm shares the shape but not the interior level (its undriven rings
float to \(r \approx 0.7\)); ``same structure'' is a claim about the two
winning routes, whose converged profiles are estimator-independent, not
about all trained arms sharing one interior value. The counterexample is
the §6 uniform pin (\(r^* = 2.0\)), which trains to \(1.3\times10^{-3}\)
with only 20/32 rings above the gate (the three inter-tap interiors
dark; §6). ``\(N = 32\) carries the measured profile'' therefore extends
from the initialization to the winning routes' solutions; a
uniform-damping design in the same box does not inherit that extension.

Two registered follow-ups bound what that means (PR-18 §18.6).
\emph{Output} participation is far narrower than gradient reach: ranking
rings by readout contribution (\(|c_j|\,\bar a_j\)) and zeroing them
cumulatively with the decoder frozen, the deployed solutions hold within
2× of their own error floor until \textasciitilde24--26 of 32 readouts
are gone --- the delivered function rides on \(N_\text{eff} \approx\)
\textbf{6--8 rings} (median 6 for PAT, 7 for SPSA; a registered
head-refit row recovers none of it) --- a number §7.2 must and does
carry. And the interior's contribution to the trained solution is real
but thin: a taps-only control (only the four driven rings'
\(\{\delta, \kappa_\text{ext}\}\) trainable, all else pinned at init)
reaches within 12\% of the full partition at eval-F, the full partition
better by a paired CI of \([+0.25, +1.55]\times10^{-4}\) excluding zero
--- the pre-registered ``discovers'' reading, earned by a modest margin.
The tap damping carries most of the solution, and the restricted arm is
not the full one truncated: its taps compensate asymmetrically
(\(r \approx \{1.52, 1.30,
0.74, 1.34\}\), seed-consistent, vs the full arm's near-uniform
\(\approx 1.3\)). The coarse floor, for the record, \emph{inverts} this
control's ordering (taps-only reads two symbols better at 3,840; eval-F
resolves the true one) --- §5.1's floor hazard, illustrated once more.

\subsubsection{5.8 Secondary diagnostic (registered,
deferred)}\label{secondary-diagnostic-registered-deferred}

PR-14 --- the bias/variance decomposition of each estimator's gradient
against the BPTT reference --- is registered as a secondary diagnostic
only and was not run in the core bake-off; it is a Stage-0.5-full row.
The reason it is secondary is structural: gradient-direction agreement
flatters the exact methods (adjoint, RHEL) and penalizes SPSA, whose
per-step alignment is poor by construction while its \emph{averaged}
trajectory converges (§5.3) --- scoring on cosine alone would have
reproduced the known failure mode of ranking estimators by a proxy the
task does not pay for. The fragments that exist (the adjoint's
0.994/0.925 cell-dependent cosine, §5.6; RHEL's non-dissipative-limit
recovery, §5.4) are reported where they carry mechanistic weight, and
the full decomposition belongs in an appendix when the S0.5-full rows
run.

\subsection{6. Choosing the damping operating
point}\label{choosing-the-damping-operating-point}

The D-LinOSS line's central observation --- that damping in an
oscillatory SSM is a \emph{performance knob}, not a defect to minimize
--- has a sharp physical meaning here: per-ring damping is the net loss
\(\kappa_\text{net}(\kappa_\text{ext})\), and the tunable coupler that
sets it is already in the trained partition. The registered disposition
(R-ii) therefore frames the question not as ``which damping value do we
freeze?'' but as ``does training \emph{find} the right damping within
the feasible box?'' --- and the experiment separates the two readings
with two arms, six damping points, eight seeds each, at the headline
cell under the convergence-controlled protocol of §5.1.

\textbf{The damping value matters enormously.} With
\(\kappa_\text{ext}\) \emph{pinned} (only detunings and inter-ring
couplings training), the converged error spans a factor of
\textasciitilde300 across the feasible box: SER 0.39 at light damping
(\(r = 0.2\) --- not chance, which is 0.75 on 4-PAM, but two of every
five symbols wrong) falling to \(1.3\times10^{-3}\) at the optimum ---
which sits at \textbf{deep overcoupling} (\(r^* = 2.0\); measured
operating \(\kappa_\text{net} = 3.7\kappa_i\) at the converged
solutions, the fixed-gain estimate being \(4.1\kappa_i\); about 6
samples of memory at 2 GS/s). The direction is instructive: the
equalization task needs only a 7-tap span, and the long memory the
light-damping regime supplies (30--45 samples) is actively harmful ---
stale symbols interfere. ``More memory'' is not free performance in a
dissipative recurrence; damping tunes memory \emph{to the task}, which
is precisely the D-LinOSS thesis in physical units, with the optimum on
the heavily-damped side for this task class.

\textbf{Training absorbs the knob.} In the second arm the full partition
trains inside progressively wider boxes \([r_\text{min}, r_\text{hi}]\).
Every box containing the pinned optimum reaches within the
pre-registered margin of it --- in fact reaching \(5\times10^{-4}\), the
§5 ceiling, \emph{better} than the best uniform pin: per-ring trainable
coupling finds a heterogeneous damping profile no single design value
can express. Boxes that exclude the good regime fail exactly as they
must (training cannot find what the clamp forbids). The registered R-ii
test is therefore \textbf{confirmed}: the designer's damping obligation
is to make the feasible box \emph{contain} the good regime; the
operating point itself is the trained substrate's job.

\textbf{Why the pin loses: controllability at the solutions (PR-18).} A
pre-registered diagnostic (PR-18; S0.11) reproduced both arms' converged
states bit-identically (16/16 seed-runs, eval-trace-exact) and re-ran
the S0.4-0 controllability gate \emph{at the solutions} rather than at
θ₀. The uniform pin pays for its damping in gradient reach: at
\(r^* = 2.0\) only \textbf{20 of 32} rings keep gradients above the
registered \(10^{-3}\)-of-max gate --- the three inter-tap interior
segments (rings 6--9, 15--18, 24--27) go gradient-dark (worst ratio
\(5\times10^{-7}\)), and the trained couplings do not rescue them
(converged \(\mu\) median exactly at its \(0.3\kappa_i\) init; no link
grows past \(0.35\kappa_i\)). The boxed winner escapes the trade: every
seed converges to the \emph{same} structure --- the four \textbf{driven}
rings damped hard (\(r \approx 1.3\)) and the 28 undriven rings held
light (\(r \approx 0.7\), operating
\(\kappa_\text{net} \approx 1.5\kappa_i\)) --- which keeps \textbf{30 of
32} rings above the gate (worst \(2.3\times10^{-4}\); the two residual
dark rings sit at segment midpoints) while still damping where the input
lands. The heterogeneous profile above is therefore not an
overparameterization curiosity: it is the mechanism by which the trained
substrate buys task-optimal damping \emph{and} its own trainability at
once --- a combination a uniform design value structurally cannot
express, and a second, sharper reading of the ×2.5 pinned-vs-trained
gap. The structure is not an artifact of this arm: the bake-off's
winning routes converge to the same tap-heavy profile at their solutions
(§5.7), and a registered taps-only control bounds how much of the
performance the structure's interior carries --- real but thin (§5.7):
``discovers'' survives its own falsifier, narrowly. (Settled-amplitude
participation decouples from gradient reach at the solutions --- winner
counts \(\{8, 15, 22\}\) of 32 above \(\{10^{-1}, 10^{-2}, 10^{-3}\}\)
vs \(\{4, 26, 32\}\) at θ₀ --- the trainability-relevant quantity is the
gradient gate; both are reported in S0.11.)

Two honest footnotes. The slow mid-grid configurations
(\(r = 0.3\)--\(0.5\)) had not fully plateaued at the training ceiling,
so the ×300 spread is a budget-bounded statement --- but both endpoints
and all winning configurations converged, and the R-ii verdict uses
converged points only. And the result recolors §5 slightly: a large
share of what in-situ training accomplished in the bake-off \emph{is}
finding the damping operating point (the θ₀ hold value, pinned, yields
0.038 --- seventy-seven times worse than the trained substrate). Since
the offline baseline finds \(r^*\) on its calibrated model just as well
(§5.5), the damping result strengthens the trainability story without
moving the advantage question. It also recolors §2's motivation: the
optimum sits deep in the heavily-damped regime, far from the
weakly-damped near-conservative corner where the oscillatory (LinOSS)
parameterization is most distinctive relative to plain diagonal SSMs ---
on this task class the substrate's results are evidence about the
broader dissipative diagonal-SSM class that §2.2 deliberately claims,
with LinOSS as its boundary case, not evidence for LinOSS-specific
expressivity.

\textbf{Does the optimum track task memory? A registered prediction,
failed.} Because ``excess memory is harmful'' measured on one 7-tap task
is close to tautological, we froze a transfer test (PR-17 §17.4--17.5):
a second task, T-A-L --- the same channel plus a −6 dB replica of its
past-tap profile delayed by 7 symbols, doubling the memory span to 14
--- with the registered prediction that the pinned-damping optimum moves
to \emph{lighter} damping (\(r^*_L < 2.0\)), and a frozen
30\%-separation rule for upgrading the claim. \textbf{The prediction
failed.} The T-A-L optimum does not move toward lighter damping at all:
the curve keeps falling to the heavy edge of the grid
(\(2.2\times10^{-3}\) at \(r = 3\) vs \(2.5\times10^{-3}\) at \(r = 2\),
eval-F, within 13\% --- the separation rule does not fire; Fig. F6,
dashed). What doubling the task span actually did was raise the error
floor everywhere (the harder channel) while leaving the optimal damping
regime where it was. The honest reading: on this substrate and task
class the heavy-damping optimum is \emph{robust} to a ×2 change in task
memory span --- the optimum is set by the bandwidth/interference trade
of the equalization family, not by naive span-matching --- and the
``damping tunes memory to the task'' sentence above must be read at that
class level, not as a per-task tracking law. A task family engineered to
\emph{need} long coherent memory (rather than a longer ISI to cancel)
remains the right probe.

That probe has since run (PR-19; S0.12, post-assembly): a
spread-spectrum task whose decisions integrate 31 chips --- a span far
beyond the 8-lag head's reach --- with the registered prediction that
the optimum moves light (\(r^*_D \leq 1.0\)). The prediction
\textbf{failed by degeneracy}: at the frozen SNR the task's integration
gain leaves every grid point at zero error --- even \(r = 3\), whose
\textasciitilde5-sample memory the mechanism said should be fatal,
clears on partial-correlation margin --- so the separation rule could
not fire (ledger §19.6b). The floor bounds the structural reading too:
at zero error the gradients vanish early and freeze the parameters where
they stand, and that is what is observed --- on T-D every arm ends
essentially at its initialization (taps \(\approx 0.34\), interior
\(\approx 0.30\); eight seeds plus the pilot), where T-A drove its taps
to \(\approx 1.3\). The contrast says the §5.7 profile is induced by the
task only in the weak sense that a task solved at initialization induces
nothing; it is not a second, independent instance of structure
discovery. The damping-tracks-memory hypothesis accordingly stands at
zero for two genuine tests --- one informative null (T-A-L) and one
degenerate attempt (T-D), with T-A the observation that generated it,
not a test of it. We leave it as a hypothesis one probe has declined to
confirm and a second could not reach.

\subsection{7. Does it pay? The systems
envelope}\label{does-it-pay-the-systems-envelope}

\subsubsection{7.1 The question, and how we keep it
honest}\label{the-question-and-how-we-keep-it-honest}

A photonic SSM only matters if, after paying the conversion toll --- DAC
and modulator in, photodiode and ADC out, heaters held all the while ---
it still beats a competent digital implementation on the axis the niche
cares about. We price this at two frozen corners (OPT = best published
device class; CONS = named vendor parts at ENOB-at-speed, never nominal
bits), against named baselines (Microsoft Brainwave's author-stated
batch-1 streaming efficiency; a coherent-DSP ASIC class; Jetson AGX
Orin), with every number traced to a frozen source row and the five
known exclusions (laser wall-plug, the substrate's Er:Si₃N₄ gain pump,
locking, control compute, packaging) explicitly unbudgeted --- they only
shrink positive cells, so negative findings are robust to them. An
assembly-time retrieval bounds their magnitude from primary sources
(supplementary ledger): the integrated-class stack --- hybrid-laser
wall-plug 0.2--0.6 W at mW-class on-chip power, TEC hold 0.18 W,
microcontroller-class control 0.13 W, FPGA-class locking up to
\textasciitilde10 W --- totals of order 0.5--1 W, the \emph{same order
as the budgeted N = 128 photonic power itself} (\textasciitilde0.9 W at
2 GS/s). The fifth item --- the 1480/980 nm pump the substrate's g = 0.9
κᵢ erbium stage needs at every ring --- was missing from the original
list and was added at the Stage-0b re-envelope (supplementary ledger;
0.8--17 mW electrical per ring near transparency, derived): 26--550 mW
at N = 32, the same order again, and larger at N = 128. Charging it
would compress the positive cells' margin (§7.2) toward single digits,
and a benchtop realization (40--100 W laser and instrument lock) would
erase the niche outright. The niche verdict below therefore carries an
integrated-realization condition on all five excluded items, alongside
the heater-class condition it already states.

\subsubsection{7.2 Inference: a conditional niche, gated by the heater
class}\label{inference-a-conditional-niche-gated-by-the-heater-class}

The original lite envelope suggested \textbf{a conditional low-latency
niche}. The later function-matched inline follow-up (§7.4) closes the
registered product window; the original calculation is reported below
for provenance. At the registered scale grid, the photonic side clears
the strongest streaming baseline (Brainwave) by up to \textasciitilde15×
in energy per sample at N = 128 and 2 GS/s, with an optical latency
lower-bound estimate of tens of nanoseconds per sample, compared with
the baseline's reported milliseconds; this is not a measured end-to-end
latency advantage --- but only when three conditions hold
simultaneously: line rates ≳0.5 GS/s (every scenario loses everything at
0.1 GS/s), N ≳ 32 (conversion is N-independent; digital cost scales with
N --- the structural effect the architecture banks on), and
\textbf{suspended low-power heaters} (\textasciitilde1 mW/π, class B).
One measured caveat now bounds the middle condition: at the headline
cell the deployed equalizer's \emph{output} rides on
\(N_\text{eff} \approx 6\)--\(8\) of 32 rings (readout ablation on the
stored solutions, PR-18 §18.6b; head-refit recovers none of the zeroed
readouts). The digital baselines in this comparison are priced at the
nominal N, but a baseline built to the \emph{function} would carry
\textasciitilde6--8 states and shrink its cost accordingly --- so the
N-scaling premise holds only for workloads that actually exercise the
state dimension, which the 7-tap equalization family does not. The
registered long-coherent-memory follow-up (PR-19: unipolar despread-31,
run to discharge exactly this condition) did \textbf{not} discharge it
--- and produced no second measurement of concentration either. At the
frozen 28 dB the task's \textasciitilde16-chip integration gain leaves
the entire damping grid error-free (ledger §19.6a/b), and on a task
solved with that much margin the readout-ablation count measures the
margin, not the substrate's utilization: the recorded
\(N_\text{eff} = 4\) (every seed; a head-refit recovers the target from
as few as 3 rings, where on the equalization task the same refit
recovers nothing) is a statement about task slack, uninformative about
concentration in either direction. (The pre-written no-rise text,
drafted for an informative null, asserted the concentration reading; it
is withdrawn by ledger amendment §19.6c.) PR-19 therefore leaves the
premise exactly where §18.6b put it --- the one informative
\(N_\text{eff}\) measurement in this program is the 6--8 above. The
\textasciitilde15× at N = 128 was computed against a baseline priced at
nominal N; a baseline built to the function would shrink by roughly the
same factor --- \textbf{the advantage is unmeasured in magnitude, not
merely conditional} --- and the burden has flipped: demonstrating a
workload that genuinely exercises N ≳ 32 states on this substrate (e.g.
the same family at a registered harder operating point) is what would
restore the premise, and no such demonstration exists in this program's
data. Under the registered worst-case holding convention the
foundry-standard heater class loses to every baseline everywhere in the
window at the deployable corner; the demonstrated-foundry path therefore
does not reach the energy niche as computed --- a Stage-1 platform
constraint stated as such, with the expected-value-holding sensitivity
(under which the optimistic corner clears in-window) reported alongside
per the review finding. Jetson's \emph{peak} rating is never beaten
anywhere; the niche claim rests on measured sustained behavior, for
which batch-1 recurrent workloads on edge GPUs are documented at
\textgreater100× below claimed peak --- page-verified: measured batch-1
GRU throughput of 1.9 and 3.5 GOp/s against claimed peaks of 0.5 and 0.8
TOp/s on the two Jetson-class devices (ratios ≈263× and ≈229×), with the
source's own conclusion stating ``a factor of over 100X'' {[}34{]}; the
Orin DLA path falls back to GPU for recurrent layers. One boundary cell
(the DSP-class comparison at N = 32) clears by \textasciitilde1\% and is
treated as a tie.

\subsubsection{7.3 Training: partial energy budgets, unresolved
wall-clock
cost}\label{training-partial-energy-budgets-unresolved-wall-clock-cost}

The device-pass and digital ledgers yield the following component
estimates at the optimistic corner. These are not total training
energies.

{\def\LTcaptype{none} % do not increment counter
\begin{longtable}[]{@{}lll@{}}
\toprule\noalign{}
route & conversion energy & digital-twin compute estimate \\
\midrule\noalign{}
\endhead
\bottomrule\noalign{}
\endlastfoot
PAT & 0.57 mJ & ≈13--44 J \\
adjoint & 1.1 mJ & 0 in the idealized physical-adjoint ledger \\
SPSA & 2.6 mJ & no digital twin \\
\end{longtable}
}

PAT's digital estimate uses the registered FLOP model and named
accelerator-efficiency classes; it is not a measured implementation. The
physical-adjoint row remains conditional on realizing the reverse pass.
SPSA's vendor-part conversion estimate is 99 mJ.

SPSA's 176,000 passes, each containing 256 samples at 2 GS/s, occupy
\textbf{22.528 ms of optical sequence time}. This is a lower bound on
training duration, not elapsed wall time. The earlier wording treated it
as wall time and inferred a total of tens of millijoules; that
total-energy claim is withdrawn (supplementary N8). A complete estimate
must include parameter writes and actuator settling between
perturbations, measurement and controller latency, reset gaps, and the
laser, gain-pump, locking, packaging, and thermal-hold power over the
full duration. In symbols,
\(E_{\rm train}=E_{\rm conversion}+E_{\rm digital}+E_{\rm writes}
+\int_0^{t_{\rm wall}}P_{\rm hold}(t)\,dt\), without double-counting
components. Those timing and write-energy quantities have not been
established for a device.

The component ledger makes PAT's digital-twin cost visible and motivates
measuring SPSA's full hardware loop. It does not establish a
four-orders-of-magnitude advantage in total training energy. RHEL's
registered pump adder is likewise a component estimate; its accuracy
censoring is unchanged.

\subsubsection{7.4 Registered inline follow-up: no energy-advantage
window}\label{registered-inline-follow-up-no-energy-advantage-window}

The later Stage 0b envelope (PR-20, frozen at \texttt{ee137c4} before
calculation; results \texttt{8931d6c}) replaces the block-level
comparison with the per-tap digital equalizer an inline device would
replace. It charges the registered actuator, common-mode
thermal-management, gain-pump, and maintenance rows. This is a
calculated model envelope, not measured hardware performance.

Across the registered 0.1--2 GS/s grid, the best optimistic class-C cell
has \(E_{\rm digital}/E_{\rm photonic}=0.64\): 3.2 pJ/sample digital
versus approximately 5.0 pJ/sample photonic at 2 GS/s, 64 taps, passive
C-3, and 10 mW global thermal hold. Thus the photonic estimate is
\textbf{1.56 times the digital energy}, missing both parity and the
registered 3-times advantage threshold. Removing the unsourced
optimistic rows lowers the best ratio to 0.23 (about 4.35 times the
digital energy). The conservative product-window maximum is 0.38. The
PR-20 kill gate fires; S0b.1--S0b.3 were not run, and no
product-simulation or fabrication spend followed. The separate S0.13
correction rerun repairs P1 evidence; it does not reopen Stage 0b. The
full ledger and results accompany the paper in \texttt{docs/s0b/} and
\texttt{results/s0b\_0/}.

Thermal management dominates the most favorable cell even with zero
actuator hold. The frozen minimum-damping model gives a
pumped-to-passive memory ratio of \((1+2r_{\min})/(0.1+2r_{\min})=3.14\)
at \(r_{\min}=0.1606\); removing gain does not give a tenfold reduction
at this operating point. These corrections change interpretation, not
the frozen arithmetic or gate.

The maintenance term inherits §7.3's optical-time lower bound.
Accounting for longer control and settling time can only increase
modeled photonic energy at a fixed retraining cadence, so it cannot
reverse this negative gate. It does mean that the calculation has not
established that practical retraining is cheap or that the chosen
cadence maintains accuracy.

For fixed power and fixed workload tap count the energy ratio grows
linearly with sample rate. Quadratic scaling is an approximation only
while usable tap count also grows proportionally with rate; finite
tap-grid limits and task capacity interrupt that scaling. Higher-rate
scenarios are outside PR-20 and have no trainability result. They remain
dormant pending a concrete workload and collaborator.

The earlier §7.2 nominal-dimension envelope is retained as the
historical calculation. This function-matched follow-up supersedes its
positive product interpretation within the registered inline window. The
simulation trainability result survives; energy-competitive hardware has
not been demonstrated.

\subsection{8. Limits of this model}\label{limits-of-this-model}

\subsubsection{8.1 Robustness here means robustness to what we
modelled}\label{robustness-here-means-robustness-to-what-we-modelled}

Every robustness statement in §5 is conditioned on the substrate of §3:
the methods absorb the imperfections \emph{we simulated} --- saturating
gain at a registered operating point, Langevin ASE at NF 7 dB, resolved
backscatter doublets, 5\%-class calibration mismatch, fresh noise per
pass. Hardware contains channels we did not model: thermal transients
and self-heating at operating power, polarization rotation, fabrication
disorder beyond the derate row, drift at cadences between our episode
and training scales, and mode-splitting behavior that is not captured by
a single always-on γ per process class. Any of these could reorder the
§5 ranking on a real chip. We regard the ranking as a hypothesis the
Stage-1 hardware exists to test, with PAT and SPSA committed precisely
because they are the two routes whose chip-level robustness is already
literature-established {[}3, 2{]}.

\subsubsection{8.2 Anchor risks and verification debts, by
name}\label{anchor-risks-and-verification-debts-by-name}

The substrate's realism leans on anchors with stated residual risks: the
C-2 loss class transfers a wide-multimode racetrack result to a
single-mode registry ring (priced by the ×2-loss derate row, §3.3); the
Er:Si₃N₄ noise budget rests on a \textbf{single coupling-loss-limited
measured NF (\textasciitilde7 dB)} in the flagship device paper --- debt
\#3's original ``no measured NF'' premise was found false at the S0.3-0
recon and our NF-A = 7.0 dB was frozen to match the measurement (§3.2);
the narrower residue (intrinsic amplifier NF not isolated) keeps the NF
sensitivity rows registered; the recurrent adjoint pass is charged
\emph{as if realizable} with no demonstration in the literature (debt
\#4 --- an inference from absence, time-stamped mid-2026); and the
white-space claim itself is one-sided evidence from a pre-registered
search, to be re-swept before submission (debt \#1). The S0.4-0
calibration retired one internal debt (the drive build-up controversy
resolved by measurement: ×0.42, doublet-quenched) and left one open: the
on-resonance floor calibration under worst-case de-saturation
(anchor-risk vii, §3.6), carried as a label --- and later made concrete
by the PR-18 endpoint diagnostic: three of eight SPSA solutions end with
a ring in the hypothetically super-threshold corner (min
\(-0.06\,\kappa_i\); §3.6), an endpoint-only measurement whose Stage-1
mitigation is the already-computed δ-aware clamp (\(r \approx 0.2547\)).

\subsubsection{8.3 The benchmark anchor we do not
use}\label{the-benchmark-anchor-we-do-not-use}

An early gate required reproducing a published LinOSS benchmark as an
external anchor. Running the authors' own code on their published seeds
reproduced the Heartbeat headline within its band; on the long-sequence
EigenWorms task the rerun mean itself lands within one published
standard deviation (\(90.6\) vs \(95.0 \pm 4.4\)) --- what fails is not
the mean but the \emph{dispersion} (per-seed \(\sigma = 9.3\),
\(2.1\times\) the published value, from a bimodal seed population) and
the mechanism behind it: a numerical-precision failure mode in the
training loss (an absorbing zero-gradient state in fp32) that makes the
published number seed-unstable rather than unreproducible (Fig. S1; a
separate reproducibility note is in preparation). The gate was
adjudicated purpose-served-with-anchor-void: all downstream accuracy
references in this program are therefore \textbf{in-house
BPTT-on-substrate ceilings} measured under our own protocol (§5.1),
never transferred published numbers. We flag fp32-sensitivity generally:
our substrate runs float64, and the eval-floor granularity of §5 is
symbol-count-limited, not precision-limited.

\subsubsection{8.4 Review independence}\label{review-independence}

Through 2026-07-06 every freeze in the ledger passed adversarial review
by an independent reviewer session reporting to the PI, and several
results in this paper exist because that review forced them (the
multi-tap input map, the mismatch decomposition, the
readout-differential rule). From 2026-07-07 the program ran in a
single-session mode in which the same agent performed both execution and
review, under standing PI delegation; every artifact from that period is
so labelled in the ledger, and the S0.4b/c/S0.5 findings --- including
the two honest nulls (the offline tie; RHEL's failure) --- should be
read with that reduced independence in mind. The pre-registration
discipline (thresholds frozen and committed before runs) is the
structural mitigation --- with the caveat stated plainly: from the
single-session date onward the registrar and the registrant are the same
agent, so the commit trail is self-graded until it is externally
anchored. That anchor was treated as a precondition, not an afterthought
--- a timestamp that follows disclosure certifies nothing --- and it is
in place in two parts: the ledger head was anchored by an OpenTimestamps
proof (\texttt{timestamps/head\_2026-08-05.txt.ots}, committed before
the external review recorded here), which certifies existence-by-date
independently of who can read the repository; and the repository itself,
including the full ledger and its complete commit history, is made
public at \texttt{github.com/LTalandier/Project\_SSM} at submission (it
was first published 2026-08-02 and withdrawn to private on 2026-08-04
for the PI's content review, which is why the date of record is the
submission date rather than the earlier one). Public release will make
the recorded commit ordering inspectable. The timestamp certifies
existence of the anchored material; it does not independently establish
the actual execution times of experiments. The post-assembly external
review recorded in the ledger (seven rounds, 2026-08-01 through
2026-08-17) doubles as a measurement of this structure. Its first five
rounds found six errata --- control-cell numbers transplanted into
headline contexts, a cross-scope ratio, a protocol-mixing claim --- all
in unregistered connective prose, none touching a pre-registered number,
rule, or verdict. The seventh found the failure mode that diagnosis
could not exclude: a pre-written consumption text (PR-19's no-rise
branch), drafted for an informative outcome, was applied verbatim to a
floor-degenerate one and presented an uninformative \(N_\text{eff}\) as
confirming evidence --- an error inside registered prose (though not
inside any number, rule, or verdict), and one that overstated the
paper's \emph{caveat} rather than its claim, which is why it survived
six rounds; it is withdrawn by ledger amendment (§19.6c), with two
same-root instances in connective prose and stale values inside rendered
figure annotations corrected alongside. The discipline held where it was
applied; its seams --- connective prose, outcome branches its texts did
not anticipate, and derived artifacts --- are now named from measurement
rather than assumed absent.

\subsubsection{8.5 Scope limits we chose}\label{scope-limits-we-chose}

The task family is deliberately narrow (continuous-signal channel
equalization plus a synthetic memory family; the registered secondary
task is deferred), and the C-3 128-ring cell never gates anything. A
registered long-coherent-memory follow-up (PR-19, despread-31) ran
post-assembly and is consumed in §6/§7.2; its chief surviving lesson is
methodological --- a follow-up task's operating point must be registered
against its own processing gain (the frozen 28 dB left the entire
damping grid error-free) --- and the harder-operating-point variant
remains an open registration, not a claim. Two axes the bake-off itself
held fixed were measured afterward in pre-registered follow-ups (§5.5):
calibration-mismatch sensitivity (corrected after the command-binding
defect was found, using the same frozen rule; supplementary N8) and
drift (a literature-calibrated random walk under a deploy-then-drift
protocol, two correlation regimes). Drift remains unmodelled
\emph{during} training at the bake-off cadence, and the tested drift
magnitude is gentle (\(\approx 1.4\,\kappa_i\) accumulated) rather than
worst-case. The systems-advantage question --- whether any of this pays
once conversion overhead is counted --- is §7's; the strongest current
evidence is §5.5's mechanism triple: no advantage at any of five tested
calibration levels spanning 5--30\%, none under common-mode drift, and a
declared pre-registered \(2.42\times\) advantage specific to
uncorrelated per-ring drift --- whose real-hardware relevance rests
entirely on how uncorrelated actual on-chip drift is, an unmeasured
quantity we elevate to the sharpest Stage-1 experiment (§9). We consider
stating that plainly to be the paper's job.

The subsequent repository audit (2026-09-13; supplementary N8) found an
implementation defect in the higher-mismatch sweep and an optical-time
versus wall-time error in the energy interpretation. The invalid
higher-mismatch comparisons were withdrawn and replaced by a bounded,
versioned correction rerun (S0.13); all five fine intervals still
include zero. Total-training-energy claims remain withdrawn. This audit
and correction were implemented and checked in one session, without
independent review.

\subsection{9. Outlook: Stage 1}\label{outlook-stage-1}

\subsubsection{9.1 Hardware development
paused}\label{hardware-development-paused}

The PR-20 inline-envelope kill gate (§7.4) closes the registered product
route. No wafer or further product-development spend is planned. The
following architecture is a possible collaborator-led research
demonstrator, not an approved fabrication plan. A higher-rate study
remains dormant until a concrete workload and collaborator justify a new
registration.

The bake-off fixes the Stage-1 chip's training stack by evidence rather
than taste: \textbf{PAT and SPSA, nothing else in the loop.} Neither
exotic route earned promotion (§5.3), and the hardware ledger
(supplementary note N2) shows why that is unlikely to reverse on
hardware grounds alone: the adjoint adds circulators, phase-coherent
reverse injection, and an unsolved separation of the counter-propagating
field from the backscatter doublet; RHEL adds a pumped conjugator bank
whose power budget (≈9.6 W at 32 rings) exceeds the entire rest of the
system. SPSA's row is the quiet asset --- zero added components, zero
model burden --- so the minimal viable demonstration is: the §3 plant (N
= 8--32 rings, foundry-floor Q suffices per C-1's gate), thermo-optic
\{δ, κ\_ext, μ\} actuation, one drop-port readout chain, the four-tap
drive map of §3.6, and SPSA as the first-light training route with PAT
layered on once the twin is characterized to the 5\%-class the mismatch
protocol assumed. One operating rule is fixed by measurement in advance
rather than discovered on hardware: SPSA is the route whose converged
solutions walk rings toward the hypothetically super-threshold corner
(three of eight seeds; §3.6), so the first-light SPSA runs under the
δ-aware clamp \(r \geq 0.2547\) --- the PR-18 endpoint diagnostic
converted anchor-risk (vii) from a limitations label into this design
input.

The multi-project-wafer path is concrete: the registered cells were
chosen to be foundry-realizable (C-1 at generic-foundry loss; C-2
bounded by a demonstrated MPW result {[}20{]}), and the actuation map of
§2 uses only standard thermo-optic tuners. The E/O overhead the four-tap
drive and the eval protocol add is exactly what §7's envelope prices.

\subsubsection{9.2 What would change our
mind}\label{what-would-change-our-mind}

Three pre-registered forks, with their triggers on record: (i)
\textbf{adjoint promotion} --- if a Stage-1-adjacent demonstration
retires debt \#4 (a physical recurrent reverse pass), the PR-9 criterion
re-opens with the S0.5 data as prior; the sim says it would arrive at
ceiling-grade accuracy at 2× PAT's device cost, zero digital. (ii)
\textbf{The in-situ advantage} --- the offline absence of a resolved
difference (§5.5) across five 5--30\%-class mismatch levels sets the
burden: in-situ training earns its place on hardware only if a measured
differential justifies it: independent drift is one modeled candidate;
mismatch beyond the corrected grid and total training-energy comparisons
remain unresolved. The Stage-1 experiment should be \emph{designed to
measure exactly this differential} --- same chip, offline-deploy vs
PAT/SPSA arms --- rather than assume it. (iii) \textbf{RHEL} --- nothing
on SiN; the sim verdict (dissipation-fatal at the operating point even
with a perfect conjugator) would need a \emph{conservative} platform
regime, not a better conjugator, to reopen.

\subsubsection{9.3 Beyond the linear unselective
core}\label{beyond-the-linear-unselective-core}

The frozen architecture is deliberately the unselective LTI core ---
poles and couplings, the part photonics builds natively. The selectivity
axis (input-dependent dynamics in the Mamba direction {[}35{]}) maps
onto the same lattice as input-dependent \(C\) then \(B\) actuation and
is scoped for a later stage only behind its own gate (per-step tuning
without per-state DACs); nothing in this paper's claims depends on it.
Likewise the damping operating point (§6) and the D-LinOSS accuracy
question ride the \emph{trainable} κ\_ext axis established here rather
than new hardware.

\subsubsection{9.4 Closing}\label{closing}

The program set out to answer a narrow question with unusual
bookkeeping: can the physics of a dissipative photonic recurrence be
trained through itself, and at what honest cost? In simulation, under
pre-registered thresholds: yes --- by the two methods a chip can already
run, at device-pass costs now quantified, with the exotic routes priced
out by data and the independent-drift advantage bounded by its
simulation assumptions. The function-matched inline envelope is negative
within the registered window. A chip would test physical trainability,
but these results do not justify product development or fabrication
spend.

\subsection{Figures}\label{figures}

\begin{figure}
\centering
\pandocbounded{\includegraphics[keepaspectratio,alt={F1}]{figures/F1_architecture_pole_region.png}}
\caption{F1}
\end{figure}

\textbf{Figure F1 --- Architecture and realizable pole region.} (a) The
N = 32 coupled-ring chain (C-2 cell): trained parameter set \{δ\_j,
κ\_ext,j, μ\_j,j+1\} (ring detunings, bus couplings, inter-ring
couplings), with the resolved four-tap input map B = \{3, 12, 21, 30\}
(§3, §5.7). (b) Realizable memory (samples at 2 GS/s) versus intrinsic Q
at gain fractions g\_f ∈ \{0, 0.5, 0.9\}, the three registered cells
(C-1 foundry-floor, C-2 headline, C-3 aspirational), and the 7-tap
task-span line --- why the operating point carries gain: at g\_f = 0 the
foundry-floor cell sits at the task-span line; at the registered g\_f =
0.9 all three cells clear it with margin.

\begin{figure}
\centering
\pandocbounded{\includegraphics[keepaspectratio,alt={F2}]{figures/F2_substrate_clamp.png}}
\caption{F2}
\end{figure}

\textbf{Figure F2 --- The dissipative substrate's operating map.} (a)
Settled saturating net loss κ\_net(r) versus the fixed-gain plane, the
lasing crossing r*, the registered clamp band (r\_min = 0.1606 with
margin m\_κ = 0.05, Δr = 0.02), and the initialization point θ₀. (b)
Off-resonance de-saturation at r\_min versus θ₀ --- the quantified
anchor-risk (vii): a detuned ring at r\_min reaches κ\_net = −0.48 κ\_i,
which is why the clamp is referenced on-resonance (§3.6).

\begin{figure}
\centering
\pandocbounded{\includegraphics[keepaspectratio,alt={F3}]{figures/F3_participation_profile.png}}
\caption{F3}
\end{figure}

\textbf{Figure F3 --- Controllability is a first-class constraint.}
Per-ring gradient magnitude (relative to the maximum ring) under a
single input tap versus the resolved four-tap map. The single-drive
profile collapses geometrically with distance from the drive (ring 32
sits \textasciitilde25 orders below the maximum; participation \{1, 3,
5\}/32 rings above \{10⁻¹, 10⁻², 10⁻³\}); the resolved B = \{3, 12, 21,
30\} puts all 32 rings above the pre-registered 10⁻³ gate (worst ring
1.42 × 10⁻³, min over five drive seeds).

\begin{figure}
\centering
\pandocbounded{\includegraphics[keepaspectratio,alt={F4}]{figures/F4_sample_efficiency.png}}
\caption{F4}
\end{figure}

\textbf{Figure F4 --- Sample efficiency to target (the headline).}
Median held-out SER versus physical device passes at C-2 (8 seeds;
shaded IQR): PAT, recurrent adjoint, SPSA, RHEL (honest echo, censored
at budget), the readout-only reservoir baseline, and the offline-deploy
arm (deployed pre-trained, so it starts low; head recalibration only).
Dotted lines mark the BPTT-on-substrate ceiling (device passes = 0 by
the PR-7 convention, drawn as a level only) and the pre-registered
target SER = 5.65 × 10⁻³; the vertical line is the device-pass budget B
= 252,800.

\begin{figure}
\centering
\pandocbounded{\includegraphics[keepaspectratio,alt={F5}]{figures/F5_ranking.png}}
\caption{F5}
\end{figure}

\textbf{Figure F5 --- Ranking at matched device-pass cost.} Median
device passes to target (per-seed dots) with the digital-computation
side-ledger co-reported (hatched; PAT's twin backward = 505,600 digital
passes): PAT 38,400 \textless{} adjoint 73,600 \textless{} SPSA 176,000;
RHEL censored 0/8. The ranking answers the pre-registered promotion
question in the negative: neither exact method beats both workhorses.

\begin{figure}
\centering
\pandocbounded{\includegraphics[keepaspectratio,alt={F6}]{figures/F6_damping.png}}
\caption{F6}
\end{figure}

\textbf{Figure F6 --- Damping is a first-order design knob.} Final SER
versus uniform pinned overcoupling r (plateaued endpoints spanning
×302), the deep-overcoupling optimum r* = 2.0 (measured κ\_net = 3.7
κ\_i at the converged solutions, 4.1 fixed-gain estimate ---
\emph{excess} memory is harmful for this task), and the trainable-κ\_ext
box (R-ii): boxes containing r* train to the 5 × 10⁻⁴ ceiling, beating
every uniform pin --- the heterogeneous damping profile is found by
training, not designed. Dashed red: the T-A-L transfer test (14-tap
span, eval-F; PR-17) --- the registered prediction that the optimum
moves to lighter damping \emph{failed}; the harder task raises the floor
everywhere while the optimum stays in the deep-overcoupling plateau (r =
2--3 within 13\%), so the heavy-damping optimum is robust to a ×2
task-memory span (§6).

\begin{figure}
\centering
\pandocbounded{\includegraphics[keepaspectratio,alt={F7}]{figures/F7_envelope.png}}
\caption{F7}
\end{figure}

\textbf{Figure F7 --- Historical systems-envelope components.} (a)
Conversion-energy estimates (SPSA 2.6 mJ OPT / 99 mJ CONS; PAT 0.57 mJ
OPT) and PAT's 13--44 J digital-twin estimate. Writes, settling, and
full-duration holding costs are not included; the plot does not rank
total training energy (§7.3, N8). (b) The original nominal-N inference
comparison at 2 GS/s. The later function-matched inline envelope finds
no energy-advantage window under PR-20 (§7.4); this panel is historical
context.

\begin{figure}
\centering
\pandocbounded{\includegraphics[keepaspectratio,alt={F8}]{figures/F8_mismatch_drift.png}}
\caption{F8}
\end{figure}

\textbf{Figure F8 --- What breaks the offline tie (pre-registered
follow-ups, §5.5; eval-F protocol of record, PR-17).} (a) The complete
corrected 5--30\%-class calibration sweep (S0.13): eight seeds per arm
and level, 32 fresh PAT units, 48 unchanged controls, and a separate
exact m=1 reproducibility anchor. All five paired difference intervals
include zero; offline/PAT ratios span 0.994--1.053. No registered
advantage is declared. Lines are medians; dots show individual seeds.
The historical higher-mismatch PAT outputs remain invalid and are not
plotted (N8).

\begin{enumerate}
\def\labelenumi{(\alph{enumi})}
\setcounter{enumi}{1}
\tightlist
\item
  Deploy-then-drift, common-mode regime (σ\_step = 0.40 κ\_i per step on
  all detunings coherently): the offline laser re-lock absorbs the drift
  and keeps pace (ratio 1.34, CI including zero). (c) Independent
  per-ring drift: the re-lock cannot fix per-ring pole scatter; in-situ
  retraining holds near-ceiling while the re-locking offline baseline
  degrades with accumulated drift --- \textbf{time-integrated ratio
  2.42×, difference-CI {[}+0.71, +2.84{]} × 10⁻³ excluding zero: both
  conditions of the frozen decision rule (PR-16: median ratio ≥ 2 AND
  difference-CI excludes 0) hold} (the coarse-floor estimate, 1.84×, sat
  below the bar and is co-reported per the frozen both-protocols rule;
  coarse and fine score the \emph{same bit-identical trajectories} at
  two floors --- ledger §17.8; the ratio's own CI {[}1.7, 4.6{]} is
  reported in §5.5). The in-situ SPSA curve is context only (per-step
  re-convergence transient; no registered verdict involves it).
  Dashed/dotted lines: the pre-registered target and the matched-budget
  BPTT reference (31,600 updates, the arms' own training budget ---
  PR-17 §17.8; the bake-off's 12,000-update ceiling is protocol-local to
  §5.3 and is not drawn).
\end{enumerate}

\begin{figure}
\centering
\pandocbounded{\includegraphics[keepaspectratio,alt={S1}]{figures/S1_g3_anchor_dossier.png}}
\caption{S1}
\end{figure}

\textbf{Figure S1 --- The benchmark anchor we do not use (G3 dossier,
§8.3).} (a) Validation-accuracy trajectories of the \emph{official}
LinOSS-IM code on the published EigenWorms seeds (our rerun, 2026
stack): two of five seeds visibly collapse mid-training (the fp32
absorbing-zero-gradient mode). (b) Final test accuracy per seed against
the published 95.0 ± 4.4\%: rerun mean 90.56\%, σ 9.34 ≈ 2.1× the
published dispersion (per-seed 97.22 / 83.33 / 97.22 / 97.22 / 77.78).

\begin{figure}
\centering
\pandocbounded{\includegraphics[keepaspectratio,alt={S2}]{figures/S2_twin_mismatch_c1.png}}
\caption{S2}
\end{figure}

\textbf{Figure S2 --- Twin-mismatch decomposition at C-1 (§5.6;
eval-F).} Final SER (3 seeds) for PAT with a perfect twin, a
parametric-error (M-par) twin, and a structural-omission (M-struct)
twin: perfect and M-struct land identically (1.11 × 10⁻³) while the fine
floor resolves a small M-par excess (1.31 × 10⁻³, +18\% --- invisible at
the coarse floor). Mismatch channels remain ≈ 0 at this cell only (the
dropped gain channel grows to \textasciitilde8\% of gradient direction
at C-2, §5.6). The idealized-conjugator RHEL control clears the C-1
target with a thin margin (7.2 vs 7.3 × 10⁻³), localizing RHEL's C-2
failure to echo physics, not mechanics.

\begin{figure}
\centering
\pandocbounded{\includegraphics[keepaspectratio,alt={S3}]{figures/S3_echo_chain.png}}
\caption{S3}
\end{figure}

\textbf{Figure S3 --- The concrete echo sub-model (PR-11, §4).} (a)
Conjugation-chain waterfall at the frozen operating point (mechanism A,
shared spiral bank): ring-port extraction η\_ex², single-pass χ³-FWM
spiral conversion (0.3 W pump, 0.5 m), routing/insertion --- chain η\_c
= −22.4 dB per conjugation. (b) Per-cell feasibility ceilings for the
alternative mechanisms: resonant-ring loaded-Q ceiling from the state
bandwidth (mechanism B) and 10-ns off-chip transit amplitude survival
(mechanism C).

\begin{figure}
\centering
\pandocbounded{\includegraphics[keepaspectratio,alt={S5}]{figures/S5_rhel_r1_recovery.png}}
\caption{S5}
\end{figure}

\textbf{Figure S5 --- RHEL recovers its own theorem's limit (§5.4).}
Cosine between the RHEL update and the exact BPTT gradient as the
substrate is made progressively less dissipative: −0.75 at κ\_net T dt ≈
1.0, monotonically to +1.0000 at 0.03 --- exact recovery of the
non-dissipative limit. The registered C-2 operating point sits at κ\_net
T dt ≈ 8 (the C-1 control cell at ≈ 27), far beyond the anti-aligned
regime: the C-2 failure is dissipative-echo bias, not implementation
error.

\subsection{References}\label{references}

\begin{enumerate}
\def\labelenumi{\arabic{enumi}.}
\tightlist
\item
  Spall, ``Multivariate stochastic approximation using a simultaneous
  perturbation gradient approximation,'' IEEE Trans. Autom. Control
  37(3), 332--341 (1992)
\item
  Bandyopadhyay, Sludds, Krastanov, Hamerly, Harris, Bunandar,
  Streshinsky, Hochberg, Englund, ``Single-chip photonic deep neural
  network with forward-only training,'' Nat. Photonics 18, 1335--1343
  (2024). arXiv:2208.01623
\item
  Wright, Onodera, Stein, Wang, Schachter, Hu, McMahon, ``Deep physical
  neural networks trained with backpropagation,'' Nature 601, 549--555
  (2022)
\item
  Hughes, Minkov, Shi, Fan, ``Training of photonic neural networks
  through in situ backpropagation and gradient measurement,'' Optica
  5(7), 864--871 (2018)
\item
  Tanaka et al., ``Recent advances in physical reservoir computing: A
  review,'' Neural Networks 115, 100--123 (2019); Van der Sande,
  Brunner, Soriano, ``Advances in photonic reservoir computing,''
  Nanophotonics 6(3), 561--576 (2017)
\item
  Jayatilleka et al., ``Wavelength tuning and stabilization of
  microring-based filters using silicon in-resonator photoconductive
  heaters,'' Opt. Express 23(19), 25084--25097 (2015); Mak, Sacher, Xue,
  Mikkelsen, Yong, Poon, ``Automatic Resonance Alignment of High-Order
  Microring Filters,'' IEEE JQE 51(11) (2015); Milanizadeh, Aguiar,
  Melloni, Morichetti, ``Canceling Thermal Cross-Talk Effects in
  Photonic Integrated Circuits,'' JLT 37(4), 1325--1332 (2019)
\item
  Bueno, Maktoobi, Froehly, Fischer, Jacquot, Larger, Brunner,
  ``Reinforcement learning in a large-scale photonic recurrent neural
  network,'' Optica 5(6), 756--760 (2018)
\item
  Böhm, Verschaffelt, Van der Sande, ``A poor man's coherent Ising
  machine based on opto-electronic feedback systems\ldots,'' Nat.
  Commun. 10, 3538 (2019); companion: Böhm et al., Nat. Commun. 13, 5847
  (2022)
\item
  Wu et al., ``Monolithically integrated asynchronous optical recurrent
  accelerator,'' eLight 5, 7 (2025)
\item
  Ashtiani, Idjadi, Kim, ``Integrated photonic neural network with
  on-chip backpropagation training,'' Nature 651, 927--932 (2026).
  arXiv:2506.14575
\item
  Zhao et al., ``In-Situ Trained Microring-Based Neural Networks,''
  Laser Photon. Rev.~(2025), 10.1002/lpor.202501576
\item
  Wu, Ren et al., ``Time-synthetic optical neural networks with stable
  programmable gain,'' arXiv:2507.02297 (retitled from ``A scalable and
  programmable optical neural network in a time-synthetic dimension'')
\item
  Zhang, Wang, Lederman, Shastri, Prucnal et al., ``Compact,
  reconfigurable, and scalable photonic neurons by
  modulation-and-weighting microring resonators,'' eLight 6, 6 (2026).
  arXiv:2505.11369
\item
  ``In-situ optimization of an optoelectronic reservoir computer with
  digital delayed feedback,'' ACS Photonics (2025). arXiv:2502.11126
\item
  Van Assche, Masaad, Gooskens, Sackesyn, Van Kerrebrouck, Yin,
  Bienstman, ``Real-time optical signal equalization with a silicon
  photonic spatially distributed reservoir computer,'' Nat. Photonics
  20, 1062--1069 (2026), 10.1038/s41566-026-01968-2. arXiv:2503.19911
\item
  Gu, Goel, Ré, ``Efficiently Modeling Long Sequences with Structured
  State Spaces,'' ICLR 2022. arXiv:2111.00396
\item
  Gu, Gupta, Goel, Ré, ``On the Parameterization and Initialization of
  Diagonal State Space Models,'' NeurIPS 2022. arXiv:2206.11893
\item
  Rusch, Rus, ``Oscillatory State-Space Models,'' ICLR 2025 (Oral).
  arXiv:2410.03943
\item
  Boyer, Rusch, Rus, ``Learning to Dissipate Energy in Oscillatory
  State-Space Models,'' arXiv:2505.12171
\item
  Cui, Cao, Pan, Gao, Yu, Zhang, ``Compact microring resonator based on
  ultralow-loss multimode silicon nitride waveguide,'' Adv. Photonics
  Nexus 2(4), 046007 (2023)
\item
  Gupta, Gu, Berant, ``Diagonal State Spaces are as Effective as
  Structured State Spaces,'' NeurIPS 2022. arXiv:2203.14343
\item
  Haus, \emph{Waves and Fields in Optoelectronics}, Prentice-Hall (1984)
\item
  Gorodetsky, Pryamikov, Ilchenko, ``Rayleigh scattering in high-Q
  microspheres,'' JOSA B 17(6), 1051--1057 (2000); Kippenberg, Spillane,
  Vahala, ``Modal coupling in traveling-wave resonators,'' Opt. Lett.
  27(19), 1669--1671 (2002)
\item
  Liu, Huang, Wang, He, Raja, Liu, Engelsen, Kippenberg, ``High-yield,
  wafer-scale fabrication of ultralow-loss, dispersion-engineered
  silicon nitride photonic circuits,'' Nat. Commun. 12, 2236 (2021)
\item
  Rukh, Colación, Buck, Drake, ``Process-structure-property
  relationships in subtractive fabrication of silicon nitride
  microresonators for nonlinear photonics,'' arXiv:2511.02198 (2025)
\item
  Liu, Qiu, Ji, Lukashchuk, He, Riemensberger, Hafermann, Wang, Liu,
  Ronning, Kippenberg, ``A photonic integrated circuit-based
  erbium-doped amplifier,'' Science 376, 1309--1313 (2022)
\item
  Ye, Marpaung, ``Compact multi-mode silicon-nitride micro-ring
  resonator with low loss,'' Adv. Photonics 5(5), 050503 (2023)
\item
  Talandier, multi-layer photonic-equalization manuscript (in
  preparation) + chip β-track thermal-crosstalk/perturbation-tolerance
  results
\item
  Pourcel, Ernoult, ``Learning long range dependencies through time
  reversal symmetry breaking,'' arXiv:2506.05259 (2025)
\item
  López-Pastor, Marquardt, ``Self-Learning Machines Based on Hamiltonian
  Echo Backpropagation,'' Phys. Rev.~X 13, 031020 (2023).
  arXiv:2103.04992
\item
  Riemensberger, Kuznetsov, Liu, He, Wang, Kippenberg, ``A photonic
  integrated continuous-travelling-wave parametric amplifier,'' Nature
  612, 56--61 (2022); Krückel et al., ``Continuous wave-pumped
  wavelength conversion in low-loss silicon nitride waveguides,'' Opt.
  Lett. 40(6), 875--878 (2015)
\item
  this work: PR-15 two-modality search + refresh memos, supplementary
  N3; public repository \texttt{github.com/LTalandier/Project\_SSM}
\item
  Dacha, Zhao, McNulty, Bhatt, Lipson, Gaeta, ``Frequency-stable
  nanophotonic microcavities via integrated thermometry,'' Nature
  Photonics (2025). arXiv:2506.21692
\item
  Gao, Rios-Navarro, Chen, Liu, Delbruck, ``EdgeDRNN: Recurrent Neural
  Network Accelerator for Edge Inference,'' IEEE JETCAS 10(4), 419--432
  (2020). arXiv:2012.13600
\item
  Gu, Dao, ``Mamba: Linear-Time Sequence Modeling with Selective State
  Spaces,'' COLM 2024 (Outstanding Paper). arXiv:2312.00752
\end{enumerate}

\clearpage

\section{Supplementary notes}\label{supplementary-notes}

\subsection{N1 --- Pre-registration
ledger}\label{n1-pre-registration-ledger}

\texttt{shared/preregistration.md}, included verbatim among the arXiv
ancillary records.

Ancillary filename: \texttt{shared\_\_preregistration.md}. Every
PR-block carries its status history (PROPOSED → AMEND → SIGNED/FROZEN)
and the commit that froze it; the two S0.4-close addenda (sizing;
ceiling) are dated \emph{before} the runs they govern. The ledger is the
paper's §3.7 ``ledger-as-method'' object.

\subsection{N2 --- Per-method hardware
ledger}\label{n2-per-method-hardware-ledger}

\texttt{docs/s0\_4/f8\_hardware\_ledger.md}: per-route observables,
actuators, added components, and calibration burdens (SPSA simplest →
RHEL heaviest); the structural note that adjoint/RHEL cannot win
promotion-rule (b) (strictly-simpler hardware) by construction.

\subsection{N3 --- White-space search
dossier}\label{n3-white-space-search-dossier}

\texttt{docs/s0\_L/debt1\_whitespace\_search.md} (PR-15 two-modality
search + kill-criterion, 2026-06-09) +
\texttt{docs/s0\_L/whitespace\_refresh\_2026-07-12.md} (assembly
refresh: W1 clean, W0 survives with qualifiers load-bearing; five named
near-misses dispatched in §1.1; four page-level reads registered for the
final pre-submission sweep). The 2026-07-27 page-read memo and
\texttt{docs/s0\_L/source\_refresh\_2026-09-13.md} complete the dated
record; the latter checks Zhang, Ashtiani, Van Assche, Rukh, and Dacha
against primary full texts and records a fresh bounded search. These
dossiers are included as ancillary records.

\subsection{N4 --- Registration and run
provenance}\label{n4-registration-and-run-provenance}

The repository \texttt{Project\_SSM} (branch \texttt{main}) is released
at \texttt{github.com/LTalandier/Project\_SSM} at submission. The
OpenTimestamps anchor is \texttt{timestamps/head\_2026-08-05.txt.ots}.
Public access makes the recorded commit ordering inspectable; the
timestamp establishes existence of the anchored material, not
independently verified experiment execution times.

Rule being evidenced: \textbf{every threshold/spec commit predates the
run that consumes it.}

{\def\LTcaptype{none} % do not increment counter
\begin{longtable}[]{@{}
  >{\raggedright\arraybackslash}p{(\linewidth - 4\tabcolsep) * \real{0.3333}}
  >{\raggedright\arraybackslash}p{(\linewidth - 4\tabcolsep) * \real{0.3333}}
  >{\raggedright\arraybackslash}p{(\linewidth - 4\tabcolsep) * \real{0.3333}}@{}}
\toprule\noalign{}
\begin{minipage}[b]{\linewidth}\raggedright
object
\end{minipage} & \begin{minipage}[b]{\linewidth}\raggedright
registered (commit)
\end{minipage} & \begin{minipage}[b]{\linewidth}\raggedright
consumed/measured (commit)
\end{minipage} \\
\midrule\noalign{}
\endhead
\bottomrule\noalign{}
\endlastfoot
PR-1.1 G3 adjudication (v2 signed) & \texttt{f659558} → \texttt{408e082}
& --- (S0.2 record) \\
PR-4 substrate freeze v2 (signed) & \texttt{b1af15d} → \texttt{55746b5}
& substrate build \texttt{c0b7359}/\texttt{2d0673b} \\
S0.4 packet: PR-6 v3 / PR-7 v2 / PR-5 / PR-12 (signed by delegation) &
\texttt{497f790} → \texttt{73099d5} → \texttt{a3c0c21} →
\texttt{cde3f4a} & S0.4-0 calibration \texttt{991aef0} \\
S0.4a spec + PR-5 mismatch levels & \texttt{0b8c817} & PAT/SPSA
build+smoke \texttt{502e26e} \\
S0.4b adjoint spec (gates B1--B5) & \texttt{c85fe62} & build+smoke
\texttt{ade7795} \\
PR-11 echo sub-model (proposed → verify-1 → numeric freeze) &
\texttt{965680d} → \texttt{2fab8e9} → \texttt{5ca8d52} & RHEL build
\texttt{b77aa60}, results \texttt{eb21520} \\
PR-3 target rule + PR-8 stats + PR-9 gates & \texttt{ae16f6d} & --- \\
Sizing addendum (U\_conv, B) & \texttt{7bcbcb9} & --- \\
Ceiling addendum (SER\_target, Δ\_M3) & \texttt{4fca58d} & bake-off
\texttt{15405fd}, C-1 diagnostics \texttt{ac0f1d2} \\
S0.6 damping spec (arms, grid, C7 rule) & \texttt{5a7f28b} & results
\texttt{7808596} \\
S0.7 training-envelope accounting rules & \texttt{bd37703} & results
\texttt{dc2d561}/\texttt{fae0512} \\
PR-16 drift advantage rule (pre-drift-run) & ledger 2026-07-22/24 & S0.9
drift runs \\
PR-17 eval-F + T-A-L spec & \texttt{241204a} & S0.10 runs
\texttt{017547b} \\
PR-17 §17.7 erratum (registered pre-rerun) & \texttt{ccc4385} &
corrected rerun \texttt{017547b} \\
PR-17 §17.8 decomposition + matched-reference spec & \texttt{a8d7ab9} &
matched run \texttt{967429d} \\
PR-18 converged-operating-point diagnostics (S0.11) & \texttt{b585623} &
stage-1 \texttt{ae1073f} / stage-2 \texttt{d108d5f} \\
PR-18 §18.6 follow-ups: taps-only control + N\_eff ablation (S0.11b) &
\texttt{ec2d4e3} & verdicts \texttt{21c9bb2} \\
PR-19 T-D long-coherent-memory (S0.12) & signed \texttt{e96e76c};
addenda \texttt{b55b135}/\texttt{a4e6a70} & fleet+verdicts
\texttt{f7b395a}; §19.6c round-7 amendment \texttt{a28e2d2} (no-rise
frame withdrawn as floor-degenerate; rule/verdict stand) \\
PAT command-binding correction (S0.13; original PR-5/PR-17 thresholds
retained) & code \texttt{04c67a6}; bounded rerun protocol
\texttt{bda989f} & results \texttt{174558c}; exact m=1 anchor, all 33
units complete \\
paper section drafts (post-results) & \texttt{9e8b4c4},
\texttt{db3a8c7}, \texttt{14aedff}, \texttt{77a93ca}, \texttt{436ebbf} &
--- \\
\end{longtable}
}

\subsection{N5 --- Reproducibility
statement}\label{n5-reproducibility-statement}

All simulations are float64 PyTorch on CPU. Runs executed on a
\textbf{mixed platform set} --- local x86-64 Linux and Hetzner cpx51
(shared x86-64) cloud instances --- with identical code, identical
registered seeds, and per-unit idempotent runners
(\texttt{analysis/s0\_5\_run\_one.py},
\texttt{analysis/s0\_6\_run\_one.py}) writing one JSON per (method,
cell, seed) unit; merge/statistics stages
(\texttt{analysis/s0\_5\_bakeoff.py}) are deterministic over the run
files. Evaluation uses reserved noise streams
(\texttt{EVAL\_SEED\_BASE\ =\ 900001}) never drawn during training.
Training randomness is seeded per unit; eval SER at the quantization
floor (2 errors / 3840 symbols) is platform-stable. The test suite (150
tests at S0.5-close; 165 at S0.12) pins substrate bit-identity gates
(e.g.~adjoint pass-2 forward identity, ledger counts). Total cloud spend
for all Stage-0 compute through S0.10, invoice-exact: 183 shared-vCPU
server-hours ≈ \textbf{€70 excl. VAT (€84 incl.)}, of which ≈84 h ≈ €32
productive runs (per-phase in \texttt{shared/results\_log.md}) and ≈99 h
≈ €38 one disclosed idle-server incident. (Provider cloud prices rose
≈3.9× for instances created after 2026-06-15; earlier in-session
estimates used the old rate.) The post-assembly PR-19 follow-up (S0.12)
added ≈132 server-hours ≈ \textbf{€50 excl. VAT} (rate-exact €0.3814/h;
includes one server lost to a hang with its 12 units re-run, and an
ssh-outage delay) --- Stage-0 total ≈ 315 h ≈ \textbf{€120 excl. VAT
(≈€145 incl.)}.

The S0.13 correction used four temporary CPX62 servers and 3.006
server-hours in aggregate. At the retrieved rate, rounding each server
up to one billed hour gives approximately \textbf{€1 incl.~VAT}, plus
small IPv4 charges; this is an estimate, not an invoice. All servers and
addresses were removed after retrieval. Per-unit Python and Torch
versions, source hashes, and cleanup evidence accompany the correction.
The revised suite has 174 tests. The checksum-verified reproduction
archive now includes the corrected run records, trained states, and
exact source bundle.

\subsection{N6 --- Supplementary
figures}\label{n6-supplementary-figures}

{\def\LTcaptype{none} % do not increment counter
\begin{longtable}[]{@{}
  >{\raggedright\arraybackslash}p{(\linewidth - 4\tabcolsep) * \real{0.3333}}
  >{\raggedright\arraybackslash}p{(\linewidth - 4\tabcolsep) * \real{0.3333}}
  >{\raggedright\arraybackslash}p{(\linewidth - 4\tabcolsep) * \real{0.3333}}@{}}
\toprule\noalign{}
\begin{minipage}[b]{\linewidth}\raggedright
S-fig
\end{minipage} & \begin{minipage}[b]{\linewidth}\raggedright
content
\end{minipage} & \begin{minipage}[b]{\linewidth}\raggedright
source
\end{minipage} \\
\midrule\noalign{}
\endhead
\bottomrule\noalign{}
\endlastfoot
S1 & G3 anchor-instability dossier: official-code EigenWorms val
trajectories (published seeds, collapse events visible) + final test acc
vs published 95.0±4.4 (rerun 90.56, σ 9.34) &
\texttt{results/s0\_2/gate\_i/xcheck\_official/} \\
S2 & PAT twin-mismatch decomposition at C-1 (perfect / M-par / M-struct
≡ ceiling) + rhel-ideal control &
\texttt{results/s0\_5/bakeoff\_diag\_c1.json} \\
S3 & echo conjugation-chain waterfall (−22.4 dB, mechanism A) + per-cell
Q\_L ceilings / transit survival (mechanisms B/C) &
\texttt{results/s0\_4c/pr11\_recon\_calc.json} \\
S4 (deferred) & PR-14 gradient bias/variance & deferred (S0.5-full) ---
slot reserved \\
S5 & RHEL non-dissipative-limit recovery (R1 cosine −0.75→+1.0000 vs
κ\_net·T·dt) & \texttt{results/s0\_4c/smoke.json} \\
\end{longtable}
}

Main-figure addendum: \textbf{F8} (S0.9 mismatch+drift, 3 panels) added
2026-07-27 alongside F1--F7; generator
\texttt{analysis/make\_sfigures.py}. Namespace ruling (round-3 review):
in the paper ``F8'' names \emph{only} the figure; the hardware-realism
ledger is always cited as supplementary note N2 (round-4: supplementary
sections renamed S→N so ``S2'' names only Figure S2) (its internal
finding-ID ``F8'' stays a docs-tree label), and internal ledger
sub-labels such as PR-5 ``§F7.3'' are cited as their parent PR block
only.

\subsection{N7 --- eval-F erratum audit (PR-17
§17.7)}\label{n7-eval-f-erratum-audit-pr-17-17.7}

The eval-F re-scoring shipped with one implementation erratum, disclosed
in §5.1 and made auditable here rather than attested. \textbf{Bug:} the
fine-evaluation helper re-measured the output normalization
\(y_\text{scale}\) on the \emph{trained} device; the trained readout is
calibrated to the \emph{training-time} normalization --- head and scale
are one decoder. \textbf{Signature (before, first pass):} chance-level
fine SER on exactly the arms that move \(\kappa_\text{ext}\), with
coarse SER at ceiling --- e.g.~in-situ PAT (m=1) fine \(0.667\) vs
coarse \(7.8\times10^{-4}\); BPTT ceiling fine \(0.742\) vs coarse
\(5.2\times10^{-4}\) --- while scale-consistent arms (offline-deploy;
all drift units, whose registered per-step re-measure tracks a
near-unchanged device) were unaffected. \textbf{Fix:} \texttt{train()}
exposes the head's trained normalization (\texttt{y\_scale} ledger key);
the fine evaluation consumes it. \textbf{Gate test} (in the released
suite): evaluation at the trained scale reproduces the in-training trace
evaluation to \(<10^{-12}\). \textbf{After (corrected rerun):} in-situ
PAT (m=1) fine \(8.5\times10^{-4}\); BPTT ceiling fine
\(9.8\times10^{-4}\). \textbf{Ordering:} registration \texttt{241204a} →
first pass → erratum registered + fix committed \texttt{ccc4385}
\emph{before} the 108-unit corrected rerun → results \texttt{017547b};
first-pass outputs quarantined in the results tree
(\texttt{invalid\_scalebug/}), not overwritten. The
valid-by-construction subsets (drift units; offline arms) are identical
in both passes.

\subsection{N8 --- Repository audit erratum
(2026-09-13)}\label{n8-repository-audit-erratum-2026-09-13}

\textbf{PAT mismatch scaling.} PR-5 §E required five parametric errors
to scale with m. The twin constructor scaled intrinsic loss and
backscatter, but command binding used m=1 values for detuning offset,
external coupling, and inter-ring coupling. At m=6 the command
offsets/factors were 0.05 κ\_i, 1.05, and 0.95 instead of 0.30 κ\_i,
1.30, and 0.70. Offline deployment used the intended scaled values. The
original m\textgreater1 comparisons and the claimed tie through 30\%
were withdrawn. Historical S0.9a/S0.10a files, including
\texttt{runs\_a\_patboth}, remain unchanged in the reproduction archive
and remain invalid for those comparisons. The m=1 bake-off, diagnostics,
and separate drift results are unaffected. The corrected binding stores
each twin's scaled levels; regressions check values and command
Jacobians at m=0,1,2,6.

\textbf{Correction completed (S0.13).} Protocol \texttt{bda989f}
preceded 32 fresh PAT units at m=\{2,3,4,6\} and all eight original
seeds, plus one m=1 seed-11 anchor. Each unit used 31,600 updates, the
original hyperparameters, and PR-17 fine evaluation at the trained
normalization. All 33 units finished, with no seed replacement or
tuning. The anchor exactly reproduced all 316 coarse evaluations,
coarse/fine SER, and pass ledgers, supporting reuse of 40 unaffected
offline and eight PAT m=1 records. The anchor is not an extra
statistical replicate. Source commit, versions, input hashes, and
trained-state hashes are retained. Results are recorded at
\texttt{174558c}; this is a post-result implementation repair, not a new
blind registration.

\texttt{results/s0\_13/analysis.json} replaces only the affected
mismatch block and records the imported S0.10 drift/reference blocks by
provenance. All five fine intervals include zero; ratios range from
0.994 to 1.053. Neither coarse nor fine evaluation finds a registered
crossover. Figure F8a now shows the corrected complete grid. The fine
result resolves no difference at the tested levels; it does not prove
equivalence. The original withdrawal and defect remain part of this
audit trail.

\textbf{Energy and timing.} The reported 22.528 ms is 176,000 × 256 / (2
GS/s), the optical sequence duration. It excludes parameter writes,
settling, reset gaps, measurement/controller latency, and holding power
over those intervals. The earlier ``tens of millijoules all-in'' and
total-energy ranking are withdrawn. Figure F7 retains the conversion and
digital-twin component budgets, explicitly labeled as partial. Stage 0b
inherited this lower-bound duration for maintenance; its negative energy
gate cannot improve when nonnegative missing costs are added at fixed
cadence.

\textbf{Stage 0b interpretation.} The unchanged PR-20 result is
E\_digital/E\_photonic = 0.64 (photonic 1.56× digital energy), or 0.23
with unsourced optimistic rows removed. The implemented pumped/passive
memory ratio at r\_min=0.1606 is 3.14. Quadratic rate scaling requires
usable tap count to grow with rate; fixed-tap scaling is linear. The
higher-rate residue has no trainability evidence and remains dormant.

\textbf{Reproduction.} \texttt{paper/repro/README.md} documents the
checksum-verified archive of local result records and benchmark metric
arrays, and the commands to regenerate figures and manuscript. The
archive includes invalid historical runs for audit; their presence does
not reinstate withdrawn claims. This revision was reviewed in one
session; automated tests are not an independent scientific review.

\end{document}
